\pdfoutput=1

\documentclass{article}


% if you need to pass options to natbib, use, e.g.:
%     \PassOptionsToPackage{numbers, compress}{natbib}
% before loading neurips_2023
\PassOptionsToPackage{numbers}{natbib}

% ready for submission
\usepackage[preprint]{neurips_2023}
\usepackage{float}
\usepackage{multirow}
\usepackage{graphicx}


% to compile a preprint version, e.g., for submission to arXiv, add add the
% [preprint] option:
%     \usepackage[preprint]{neurips_2023}


% to compile a camera-ready version, add the [final] option, e.g.:
%     \usepackage[final]{neurips_2023}


% to avoid loading the natbib package, add option nonatbib:
%    \usepackage[nonatbib]{neurips_2023}

\newcommand{\at}[1]{\textcolor{blue}{AT: #1}}


\usepackage[utf8]{inputenc} % allow utf-8 input
\usepackage[T1]{fontenc}    % use 8-bit T1 fonts
\usepackage{url}            % simple URL typesetting
\usepackage{booktabs}       % professional-quality tables
\usepackage{amsfonts}       % blackboard math symbols
\usepackage{nicefrac}       % compact symbols for 1/2, etc.
\usepackage{microtype}      % microtypography
\usepackage{xcolor}         % colors
\usepackage{bm}
\usepackage{wrapfig}
\usepackage{enumitem}       % Customize lists
\usepackage[per-mode=symbol]{siunitx}        % better numbers with units

\newlist{compactitem}{itemize}{3} % 3 is max-depth
\setlist[compactitem]{label=\textbullet, nosep, leftmargin=1.5em}

% Necessary packages
\usepackage{amssymb} % Beautiful symbols
\usepackage{amsmath}


\newcommand{\loss}{\mathcal{L}}

\newcommand{\R}{\mathbb{R}} % Set of real numbers
\newcommand{\GNN}{\text{GNN}}

\definecolor{cyan_our}{RGB}{0,255,255}

\usepackage[capitalise]{cleveref}

\title{Accelerating Molecular Graph Neural Networks via Knowledge Distillation}


% The \author macro works with any number of authors. There are two commands
% used to separate the names and addresses of multiple authors: \And and \AND.
%
% Using \And between authors leaves it to LaTeX to determine where to break the
% lines. Using \AND forces a line break at that point. So, if LaTeX puts 3 of 4
% authors names on the first line, and the last on the second line, try using
% \AND instead of \And before the third author name.


\author{%
  Filip Ekström Kelvinius \thanks{Equal contribution, order determined by coin flip.}\\
  Linköping University \\
  \texttt{filip.ekstrom@liu.se}
  \And
  Dimitar Georgiev \footnotemark[1] \\
  Imperial College London\\
  \texttt{d.georgiev21@imperial.ac.uk} \\
  \AND
  Artur Petrov Toshev \footnotemark[1] \\
  Technical University of Munich \\
  \texttt{artur.toshev@tum.de} \\
  \And
  Johannes Gasteiger \\
  Google Research \\
  \texttt{johannesg@google.com} \\
}


\begin{document}


\maketitle


\begin{abstract}
Recent advances in graph neural networks (GNNs) have allowed molecular simulations with accuracy on par with conventional gold-standard methods at a fraction of the computational cost. Nonetheless, as the field has been progressing to bigger and more complex architectures, state-of-the-art GNNs have become largely prohibitive for many large-scale applications. In this paper, we, for the first time, explore the utility of knowledge distillation (KD) for accelerating molecular GNNs. To this end, we devise KD strategies that facilitate the distillation of hidden representations in directional and equivariant GNNs and evaluate their performance on the regression task of energy and force prediction. We validate our protocols across different teacher-student configurations and demonstrate that they can boost the predictive accuracy of student models without altering their architecture. We also conduct comprehensive optimization of various components of our framework, and investigate the potential of data augmentation to further enhance performance. All in all, we manage to close as much as 59\% of the gap in predictive accuracy between models like GemNet-OC and PaiNN with zero additional cost at inference.
\end{abstract}

\section{Introduction}
In the last couple of years, the field of molecular simulations has undergone a rapid paradigm shift with the advent of new, powerful computational tools based on machine learning (ML) \cite{noe2020machine, westermayr2021perspective}. At the forefront of this transformation have been recent advances in graph neural networks (GNNs), which have brought about architectures that more effectively capture geometric and structural information critical for the accurate representation of molecules and molecular systems \cite{reiser2022graph, wang2023graph}. Consequently, a multitude of GNNs have been developed, which now offer predictive performance on par with conventional gold-standard methods like density functional theory (DFT) at a fraction of the computational cost at inference time \cite{batzner_e3-equivariant_2022, gasteiger_directional_2020, gasteiger_gemnet_2021, musaelian_learning_2023}. This has, in turn, significantly accelerated the modeling of molecular properties and the simulation of diverse molecular systems.

Nonetheless, this progress - largely coinciding with the development of bigger and more complex models, has naturally come at the expense of increased complexity \cite{sriram2022towards, zitnick2022spherical}. This has gradually limited the utility of state-of-the-art GNNs for large-scale molecular simulation applications, where inference throughput (i.e., how many samples can be processed for a given time) is critical for making fast continual predictions about the evolution of a system. Therefore, addressing the trade-off between accuracy and computational demand remains essential for creating more affordable tools for molecular simulations and expanding the transformational impact of GNN models in the area.

% Knowledge distillation represents a promising approach for tackling this challenge \cite{hinton_distilling_2015, gou2021knowledge}. Yet, to this date, the majority of KD research has focused on classification problems involving grid-based data, such as images and text, with applications to regression tasks and GNNs still being limited \cite{}. And, in particular, the adoption of KD for accelerating GNNs for molecular simulations is still entirely unexplored.

Motivated by that, in this work, we, for the first time, investigate the potential of knowledge distillation (KD) in enhancing the performance and scalability of state-of-the-art GNNs for molecular simulations (Fig.~\ref{fig:overview}). 
% This is a large-scale, multi-output regression task, which is challenging to address with most conventional KD methods.
To this end, we devise custom strategies for KD on molecular GNNs, which we call \emph{node-to-node (n2n)}, \emph{edge-to-node (e2n)} and \emph{vector-to-vector (v2v)} knowledge distillation. These overcome common limitations of KD for regression tasks by facilitating the distillation of hidden representations in directional and equivariant GNNs. We evaluate the performance of our KD protocols in augmenting the training process of different student models trained to predict molecular properties like energy and forces, and demonstrate the effectiveness of our protocols across teacher-student configurations, datasets and KD settings. We show that they can substantially improve the performance of student models without modifying their architecture, allowing us to reduce the performance gap between models like GemNet-OC \cite{gasteiger_gemnet_oc_2022} and PaiNN \cite{schutt_equivariant_2021} by as much as $59.16\%$ in energy predictions and $16.7\%$ in force predictions, while fully preserving the throughput of the simpler student model.

\begin{figure}[t]
  \centering
  \includegraphics[width=\textwidth]{abstract.png}
  \caption{Using knowledge distillation, we manage to significantly boost the predictive accuracy of different student models without altering their architecture. Hence, we substantially lessen the tradeoff between speed and performance in molecular GNNs, allowing us to run faster molecular simulations while maintaining predictive accuracy.}
  \label{fig:overview}
\end{figure}

In summary, the contributions of this paper are as follows:
\begin{compactitem}
    \item We explore the utility of knowledge distillation for accelerating GNNs for molecular simulations - a large-scale, multi-output regression task, challenging to address with most current KD methods.
    \item We design custom KD strategies, which facilitate the distillation of hidden representations in directional and equivariant GNNs. We demonstrate the effectiveness of our protocols across diverse teacher-student configurations and datasets.
    \item We conduct a comprehensive analysis of different components of our KD strategies, as well as explore data augmentation techniques for further improving performance.
\end{compactitem}

\section{Background and related work}
\textbf{Molecular simulations.}
In this work, we consider molecular systems at an atomic level, i.e., $N$ atoms represented by their atomic number $\boldsymbol{z} = \{z_1, . . . , z_N \} \in \mathbb{Z}^N$ and position $\boldsymbol{X} =\{ {\boldsymbol{x}_1, \dots, \boldsymbol{x}_N}\} \in \mathbb{R}^{N\times 3}$. Given a system, we want a model that can predict the energy $E \in \mathbb{R}$ of the system, and the forces $\boldsymbol{F} \in \mathbb{R}^{N\times 3}$ acting on each atom. Both these properties are of high interest when simulating molecular systems. The energy of a system is essential for the prediction of its stability, whereas the forces are important for molecular dynamics simulations, where computed forces are combined with the equations of motion to simulate the evolution of the system over time.

\textbf{GNNs for molecular systems.}
GNNs are a suitable framework for modeling molecular systems. Each molecular system $(\boldsymbol{X}, \boldsymbol{z})$ can be represented as a mathematical graph, where the set of atoms corresponds to the nodes $\mathcal{V}$, and edges $\mathcal{E}$ created between nodes by connecting the closest neighboring atoms (typically defined by a cutoff radius and/or a maximum number of neighbors). Hence, in the context of molecular simulations, we can create a GNN that operates on atomic graphs $G = (\mathcal{V}, \mathcal{E})$ by propagating information between the atoms and the edges, and makes predictions about the energy and forces of each system in a multi-output manner - i.e., $\hat{E}, \boldsymbol{\hat{F}} = \GNN(\boldsymbol{X}, \boldsymbol{z})$.

The main problem when modeling molecules and molecular properties are the number of underlying symmetries to consider.
% , most importantly rigid transformations of the atoms.
For instance, the total energy $E$ of a system is not affected by (i.e., is \emph{invariant} to) rotations and translations of the system. However, the forces $\boldsymbol{F}$ do change as we rotate a system - i.e., they are \emph{equivariant} to rotations. Therefore, to make accurate predictions about molecular systems, it is crucial to devise models that respect these symmetries. There is now a plethora of diverse molecular GNNs that achieve that, e.g., SchNet \cite{schütt2017schnet}, DimeNet \cite{gasteiger_directional_2020, gasteiger_dimenetpp_2020}, PaiNN \cite{schutt_equivariant_2021}, GemNet \cite{gasteiger_gemnet_2021, gasteiger_gemnet_oc_2022}, NequIP \cite{batzner_e3-equivariant_2022}, and SCN \cite{zitnick2022spherical}.

% review on geometric equivariant GNNs \cite{han2022geometrically}; directional message passing

% Model comparison - Inference time vs performance - table or resources, showing the tradeoff between performance and accuracy; model selection process.

%  ... More info about types of GNNs, features, problem formulation and recent results.
 
% Recently, new, well-curated datasets have been made publicly available that accelerate model research and development. These consist of millions of samples which have been synthetically generated and evaluated with DFT. This includes OC20 \cite{chanussot_open_2021} and ... Why useful, how they have been used, recent trends and results.


\textbf{Knowledge distillation.}
% In recent years, machine learning models have become exponentially larger and more complex \cite{sevilla2022compute}. This has rendered many of them prohibitive in various large-scale applications. Because of that, significant research effort has been devoted to devising strategies for compressing and accelerating ML models \cite{cheng2018model}. 
% Knowledge distillation represents a promising approach for tackling this challenge \cite{hinton_distilling_2015, gou2021knowledge}. In KD, knowledge from a complex, more competent model is transferred to a simpler one, with the objective of reducing complexity while maintaining performance \cite{}. 
Knowledge distillation is a technique for compressing and accelerating ML models \cite{cheng2018model}, which has recently demonstrated significant potential in domains like computer vision \cite{wang2021knowledge} and natural language modeling \cite{sanh2019distilbert}. The main objective of KD is to create more efficient models by means of transferring knowledge (e.g. model parameters and activations) from large, computationally expensive, more accurate models, often referred to as teacher models, to simpler, more efficient models called student models \cite{gou2021knowledge}. Since the seminal work of Hinton \textit{et al.} \cite{hinton_distilling_2015}, the field has drastically expanded methodologically, with the development of protocols that accommodate the distillation of "deeper" knowledge, more comprehensive transformation and fusion functions, as well as more robust distillation losses \cite{gou2021knowledge, hu2022teacher}. Yet, these advances have mostly focused on classification, resulting in methods of limited utility in regression tasks. Moreover, most research in the area has been confined to grid-structured data (e.g., images, text, tabular data). Despite recent efforts to extend KD to graph data and GNNs, these have likewise only concentrated on classification tasks involving standard GNN architectures \cite{tian2023knowledge}. And, in particular, the application of KD to state-of-the-art, directional and equivariant GNN architectures, where features can include complex geometric information and span both node- and edge-level features, as well as to real-world regression problems in molecular simulations, is still unexplored.

% \begin{figure}
%   \centering
%   \fbox{\rule[-.5cm]{0cm}{5cm} \rule[-.5cm]{13.5cm}{0cm}}
%   \caption{Schematic diagram of our approach depicting our custom n2n, e2n, v2v, e2e strategies.}
%   \label{fig:kd_strategies}
% \end{figure}

\section{Knowledge distillation in molecular GNNs}
\textbf{Vanilla KD.}
The standard loss function when training molecular GNNs is a loss that combines both the energy and force prediction error as follows: 
\begin{align}
    \loss_0 = \alpha_\text{E} \loss_\text{E}(\hat{E}, E) + \alpha_\text{F} \loss_\text{F}(\hat{\boldsymbol{F}}, \boldsymbol{F}),
\end{align}
where $E$ and $\boldsymbol{F}$ are the ground-truth energy and forces, $\hat{E}$ and $\hat{\boldsymbol{F}}$ are the predictions of the model of interest, and $\loss_\text{E}$ and $\loss_\text{F}$ are some loss functions weighted by $\alpha_\text{E}, \alpha_\text{F}\in\R$. 

In KD, we augment this training process by defining an auxiliary knowledge distillation loss term $\loss_{\text{KD}}$, which is added to $\loss_0$ (with a factor $\lambda\in\R$) to derive a new training loss function $\loss$ of the form:
\begin{align}
    \loss = \loss_0 + \lambda \loss_{\text{KD}}.
\end{align}
This was originally proposed in the context of classification by leveraging the fact that the soft label predictions (i.e., the logits after softmax normalization) of a given (teacher) model carry valuable information that can complement the ground-truth labels in the training process of another (student) model \cite{hinton_distilling_2015}. Since then, this has become the standard KD approach - usually referred to as vanilla KD in the literature, which is often the foundation of new KD protocols. The main idea of this technique is to employ a KD loss $\loss_{\text{KD}}$ that enforces the student to mimic the predictions of the teacher model. This is commonly achieved by constructing a loss $\loss_{\text{KD}} = \text{KL}(z_s, z_t)$ based on the Kullback–Leibler (KL) divergence between the soft logits of the student $z_s$ and the teacher $z_t$.

However, this strategy - based on the distillation of the output of the teacher model only \cite{gou2021knowledge} - poses two significant limitations. Firstly, it is by design exclusively applicable to classification tasks, since there are no outputs analogous to logits in regression setups \cite{cheng2018model, saputra2019distilling}. This has consequently limited the utility of most KD methods for regression tasks. Secondly, this approach forces the student to emulate the final output of the teacher directly, which is often unattainable in regimes where the complexity gap between the two models is substantial, and thus detrimental to KD performance \cite{cho2019efficacy}. 

\textbf{Feature-based KD.}
To circumvent these shortcomings, we focus on feature-based KD \cite{gou2021knowledge}. This is an extension of vanilla KD \cite{romero2015fitnets}, which is concerned with the distillation of knowledge across the intermediate layers of models. This allows more lightweight models to be trained to mimic features that are easier to assimilate compared to the final output directly \cite{aguilar2020knowledge}. In this paper, we perform knowledge distillation of intermediate representations by devising a loss on selected hidden features $H_\text{s} \in U_\text{s}$ and $H_\text{t} \in U_\text{t}$ in the student and teacher models respectively, which takes the form:
\begin{align}
\label{eq:kd_our}
    \loss_{\text{KD}} = \loss_{\text{feat}}(\mathcal{M}_\text{s}(H_\text{s}), \mathcal{M}_\text{t}(H_\text{t})),
\end{align}
where $\mathcal{M}_\text{s}: U_\text{s} \mapsto  U$ and $\mathcal{M}_\text{t}: U_\text{t} \mapsto U$ are transformations that map the hidden features to a common feature space $U$, and $\loss_{\text{feat}}: U \times U \mapsto \R^+$ is some loss of choice. Possible options for the transformations $\mathcal{M}_\text{s}, \mathcal{M}_\text{t}$ include the identity transformation, linear projections and multilayer perceptron (MLP) projection heads; whereas for the distillation loss $\loss_{\text{feat}}$, typical functions are mean squared error (MSE) and mean absolute error (MAE). 

The fundamental design decision when devising a KD strategy based on the distillation of internal representations is the choice of appropriate representations (i.e., features $H_s$ and $H_t$). One needs to ensure that paired features are similar both in modeling capacity and relevance to the output. Most research on feature-based distillation on graphs has so far focused on models that only have one type of (scalar) features in single-output classification tasks \cite{zhang2020reliable, he2022compressing, yu2022sail, huo2022t2}, thereby reducing the problem to the selection of layers to pair across the student and the teacher. This is often further simplified by utilizing models of the same architecture.

\textbf{Feature-based KD in molecular GNNs.} 
Most GNNs for molecular simulations contain diverse features (scalars, vectors and/or equivariant higher-order tensors based on spherical harmonics) organized across nodes and edges within a complex graph topology. These are continually evolved by model-specific operators to infer molecular properties, such as energy and forces, in a multi-output prediction fashion. Therefore, features often represent different physical, geometric and/or topological information relevant to specific parts of the output. This significantly complicates the design of an effective KD strategy, especially when the teacher and the student differ architecturally.

In this work, we set out to devise KD strategies that are representative and effective across various molecular GNNs. We consider GNNs that have diverse architectures and performance profiles, and can be organized in teacher-student configurations at different levels of architectural disparity. This is why we investigate the effectiveness of KD with respect to the following three GNN models:
% , ordered by computational complexity (ascending):

\begin{compactitem}
\item \emph{SchNet} \cite{schutt_schnet_2017}: A simple GNN model based on continuous-filter convolutional layers, which only contains scalar node features $\boldsymbol{s}\in \R^d$. These are used to predict the energy $\hat{E}$. The force is then calculated as the negative gradient of the energy with respect to the atomic positions, i.e.,  $\hat{\boldsymbol{F}} = -\nabla \hat{E}$.
\item \emph{PaiNN} \cite{schutt_equivariant_2021}: A GNN based on equivariant message passing, which contains scalar node features $\boldsymbol{x}\in \R^{d_1}$ - used for energy prediction; as well as geometric vectorial node features $\boldsymbol{v} \in \R^{3\times d_2}$ that are equivariant to rotations and can thus be combined with the scalar features to make direct predictions of the forces (i.e., without computing gradients of the energy).
\item \emph{GemNet-OC} \cite{gasteiger_gemnet_oc_2022}: A GNN model that utilizes directional message passing between scalar node features $\boldsymbol{h}\in\R^{d_\text{h}}$ and scalar edges features $\boldsymbol{m}\in\R^{d_\text{m}}$. After each block of layers, these are processed through an output block, resulting in scalar node features $\boldsymbol{x}_{\text{E}}^{(i)}$ and edge features $\boldsymbol{x}_{\text{F}}^{(i)}$, where $i$ is the block number. The output features from each block are aggregated into output features $\boldsymbol{x}_E$ and $\boldsymbol{x}_F$, which are used to compute the energy and forces respectively.
\end{compactitem}

An overview of the features of the three models can be found in \cref{tab:featuretypes}

\begin{table}[h]
\centering
\caption{\label{tab:featuretypes} Molecular GNNs can have diverse features depending on their architecture. This is an overview of the types of features available in the three models we use in this study.}
\begin{tabular}{l c c c}
 & SchNet & PaiNN & GemNet-OC  \\
 \cline{2-4} \\
Scalar node features & \checkmark & \checkmark & \checkmark \\
Scalar edge features & & & \checkmark \\
Vectorial node features & & \checkmark & \\
Output blocks & & & \checkmark \\
\end{tabular}
\end{table}


\textbf{Defining feature distillation strategies.} 
As the three models considered in this work contain very different kinds of features, there is room for considering different ways of distilling the information contained in the respective features between model architectures. In the context of the three model architectures we have considered and their types of features, we devise three different KD strategies: 

\underline{\emph{node-to-node (n2n):}} As all three models considered in this study contain scalar node features $H_{\text{node}}$, we can distill knowledge in between these directly by defining a loss $\loss_{KD}$, such that
\begin{align}
\label{eq:n2n}
    \loss_{KD} = \loss_{\text{feat}} (\mathcal{M}_{s}(H_{\text{node}, \text{s}}), \mathcal{M}_{t}(H_{\text{node},\text{t}})).
\end{align}

\underline{\emph{edge-to-node (e2n):}} The GemNet-OC model heavily relies on its edge features, which are a key component in the directional message passing defined in the architecture and can be useful as a KD resource. However, the other models considered here do not have similar edge features to distill to. To accommodate that, we propose a KD strategy where we transfer information from GemNet-OC's edge features $H_{\text{edge}, (i,j)}$ by first aggregating them as follows:
\begin{align}
    H_{\text{edge2node}, i} = \sum_{j\in\mathcal{N}(i)} H_{\text{edge}, (i,j)},
\end{align}
where $i$ is the node index. The resulting vector $H_{\text{edge2node}, i}$ is a scalar, node-level feature, and we can, therefore, use it to transfer knowledge to the student node features $H_{\text{node}, \text{s}}$ as in \cref{eq:n2n}.

\underline{\emph{- vector-to-vector (v2v):}} Similarly, the PaiNN model defines special vectorial features, which differ substantially from the scalar (node and edge) features available in the other models. These are not scalar and invariant to rigid transformations of the atoms, but geometrical vectors that are equivariant with respect to rotations. This poses a new challenge when distilling knowledge from or onto PaiNN. To this end, we define a KD procedure, where we transfer knowledge between (equivariant) vectorial node features and (invariant) scalar edge features. We achieve that by noting that scalar edge features sit on an equivariant 3D grid since they are associated with an edge between two atoms in 3D space. Hence, we can aggregate the edge features $\{H_{\text{edge}, (i,j)}\}_{j\in \mathcal{N}}$ corresponding to a given node $i$ into node-level equivariant vectorial features $H_{\text{vec}, i}$ by considering the unit vector $\boldsymbol{u}_{ij} = \frac{1}{|\boldsymbol{x}_j - \boldsymbol{x}_i|}(\boldsymbol{x}_j - \boldsymbol{x}_i)$ that defines the direction of the edge $(i,j)$, such that
\begin{align}
    H_{\text{vec}, i}^{(k)} = \sum_{j\in\mathcal{N}(i)} \boldsymbol{u}_{i,j} H_{\text{edge}, (i,j)}^{(k)},
\end{align}
with the superscript $k$ indicating the channel. This fulfills the condition of equivariance with respect to rotations as each vector $\boldsymbol{u}_{ij}$ is equivariant to rotations, with $H_{\text{edge}, (i,j)}^{(k)}$ - a scalar not influencing its direction.

\textbf{Baseline KD strategies.}
To validate the performance of our KD strategies, we evaluate their performance against 2 vanilla-based KD approaches suitable for regression tasks.

\underline{\emph{Vanilla (1):}} As mentioned above, the main problem with using vanilla KD for regression is the lack of features analogous to logits. One way of adapting vanilla KD for regression is by steering the student to mimic the final output of the teacher directly \cite{xu2022contrastive}:
\begin{align}
    \loss_{\text{KD}} = \alpha_\text{E} \loss_\text{E}(\hat{E}_\text{s}, \hat{E}_\text{t}) +  \alpha_\text{F} \loss_\text{F}(\hat{\boldsymbol{F}}_\text{s}, \hat{\boldsymbol{F}}_\text{t}),
\end{align}
where the subscripts $_\text{s}$ and $_\text{t}$ refer to the predictions of the student and teacher, respectively. Note that, unlike in classification, this approach does not provide much additional information in regression tasks, except for some limited signal about the error distribution of the teacher model \cite{cheng2018model, saputra2019distilling}.

\underline{\emph{Vanilla (2):}}
One way to enhance the teacher signal during training is to consider the fact that many GNNs for molecular simulations make separate atom- and edge-level predictions which are consequently aggregated into a final output. For instance, the total energy $E$ of a system is usually defined as a sum of the predicted contributions from each atom $\hat{E} = \sum_i \hat{E}_i$. Hence, we note that we can extend the aforementioned vanilla KD approach by imposing a loss on these granular predictions instead. Following the energy definition above, the KD loss can be expressed as:
\begin{align}
    \loss_{\text{KD}} = \frac{1}{N} \sum_{i=1}^N \loss_\text{E}(\hat{\text{E}}_{i,\text{s}}, \hat{E}_{i,\text{t}}).
\end{align}
These individual energy contributions are not part of the labeled data, but, when injected during training, provide more fine-grained information than the aggregated prediction.

\section{Experimental results}
To evaluate our proposed methods, we perform experiments on the OC20-2M \cite{chanussot_open_2021} dataset - a representative subset of the diverse OC20 dataset for molecular simulations \cite{gasteiger_gemnet_oc_2022}. We focus on the S2EF task, i.e., predicting energies and forces for a given system (atom types $\boldsymbol{z}$ and positions $\boldsymbol{X}$), which is the task of broadest interest \cite{chanussot_open_2021}. We use the models as implemented in the OC20 codebase\footnote{\url{https://github.com/Open-Catalyst-Project/ocp}} 
% We use the default values of $\alpha_\text{E}=1$ and $\alpha_\text{F}=100$ for weighting the energy- and force losses.
(see \cref{ap:train_hyperparams} for detailed information about the training procedure and different model hyperparameters). We then validate the robustness of our approach by conducting additional experiments on the COLL dataset introduced in \cite{gasteiger_dimenetpp_2020}.

\textbf{Benchmarking molecular GNNs.} 
We selected SchNet, PaiNN and GemNet-OC for our experiments as they cover most of the accuracy-complexity spectrum. To highlight this, we start by first benchmarking the predictive accuracy and computational cost of the three models on the OC20-2M \cite{chanussot_open_2021} dataset. In conjunction with the default PaiNN model - referred to as PaiNN-big in this paper, we also utilize a more lightweight version of the architecture (referred to as PaiNN-small), where we reduce the number of hidden layers and their dimensionality. We train SchNet, PaiNN-small and PaiNN-big to convergence ourselves, whereas the GemNet-OC model we employ is the pre-trained model as available within the OC20 repository.
The benchmarking results, summarized in Table \ref{tab:baselines}, 
% show the clear tradeoff between predictive accuracy and computational cost. 
% In particular, GemNet-OC is the best-performing model with an energy MAE of 286 \si{\milli\electronvolt} and a force MAE of 25.7 \si{\milli\electronvolt\per\angstrom}. However, it is also the most computationally expensive model, yielding a throughput of about 75.8 samples per GPU second at inference. At the other end of the spectrum, we have the SchNet model as the most efficient architecture, offering a $10\times$ throughput compared to GemNet-OC at 788.2. Yet, this comes at the expense of predictive performance, with results of 1308 \si{\milli\electronvolt} and 65.1 \si{\milli\electronvolt\per\angstrom} for the energy and force errors respectively. In the middle of the spectrum, we have the default PaiNN model (PaiNN-big), which achieves an energy MAE of 440 \si{\milli\electronvolt} and a force MAE of 45.3 \si{\milli\electronvolt\per\angstrom} at a throughput of 237.8 samples per GPU second. Interestingly, despite substantially reducing the size of the default PaiNN architecture in the smaller PaiNN-small model to achieve almost $3\times$ boost in throughput, the drop in performance of the more lightweight model is marginal - around 11\% for energy and 4\% for force predictions. These results 
demonstrate the substantial tradeoff between predictive accuracy and computational cost across different GNN architectures, and motivate the need for methods that can alleviate this limitation. 


\begin{table}[t]
\centering
\caption{Evaluation of the performance of the four baseline models on the OC20 S2EF validation dataset (trained on OC20-2M), averaged across OC20's four validation sets. The results show a clear tradeoff between accuracy and computational cost. For individual validation sets, see Appendix \ref{ap:main_results_full}.}
\begin{tabular}{lcccccc}
\multicolumn{1}{c}{\multirow{4}{*}{}} & & Inference & \multicolumn{4}{c}{\multirow{2}{*}{OC20 S2EF Validation}} 
\\
\multicolumn{1}{c}{} & & Throughput & \multicolumn{4}{c}{}
\\\cmidrule(r){3-3}\cmidrule(r){4-7}
\multicolumn{1}{c}{}   &   &               Samples /      &      Energy MAE  &    Force MAE  & Force cos & EFwT  \\
  \multicolumn{1}{l}{Model}     &   &         GPU sec. $\uparrow$          &      \si{\milli\electronvolt} $\downarrow$   &    \si{\milli\electronvolt\per\angstrom} $\downarrow$     & $\uparrow$  & \% $\uparrow$  \\\cmidrule(r){1-7}
SchNet             &   &           $788.2$          &   $1308$     &   $65.1$    &  $0.204$  &  $0$ \\
PaiNN-small              &  &        $618.2$             &    $489$      &     $47.1$     &  $0.345$ &  $0.085$  \\
PaiNN-big              &   &        $237.8$             &   $440$     &   $45.3$      &  $0.376$   & $0.14$ \\
GemNet-OC            &  &           $75.8$          &    $286$      &     $25.7$   &  $0.598$    & $1.06$  
\end{tabular}
\label{tab:baselines}
\end{table}


\begin{wrapfigure}{R}{0.55\textwidth}
    \includegraphics[width=1\linewidth]{cka/baseline_models.png} 
    \caption{Similarity analysis between the node features of SchNet, PaiNN and GemNet-OC using CKA.} 
    \label{fig:cka}
\end{wrapfigure}

\textbf{Similarity analysis.} 
To make our analyses exhaustive, we set out to design experiments involving teacher and student architectures of a variable degree of architectural disparity. As a proxy of that, we derive similarity scores based on central kernel alignment (CKA) \cite{kornblith2019similarity, raghu2021do}. In particular, we calculate the CKA similarity score across the hidden node representations of the trained SchNet, PaiNN and GemNet-OC (averaged over $n=987$ nodes). The results of this analysis are summarised in \cref{fig:cka}.
% Later in this section, we also investigate how we can expand on this analysis to gain explainability into the effect of KD.
Focusing on the intra-model similarities first (the diagonal in \cref{fig:cka}), we observe that PaiNN and SchNet exhibit similar behavior, with representations having a generally high degree of similarity within the two models.
This indicates that the information captured across the respective architectures is similar, which is consistent with the fact that the features of PaiNN and SchNet are iteratively updated by adding displacement features computed at each layer. GemNet-OC, however, displays completely different behavior, where the representations extracted at each layer are significantly different from the rest of the architecture. 
This is likewise in line with the GemNet-OC model, where each layer embedding is a separate output feature. When examining the similarity between different models instead, we see that, generally speaking, the node features in SchNet and PaiNN are similar, whereas those between SchNet and GemNet-OC, and PaiNN and GemNet-OC, diverge significantly as we move deeper into GemNet-OC. 

\textbf{KD setup.}
Based on the aforementioned analyses, we define the following teacher-student pairs, covering the whole spectrum of architectural disparity: PaiNN-big to PaiNN-small (\emph{same} architecture); PaiNN-big to SchNet (\emph{similar} architectures); GemNet-OC to PaiNN-big (\emph{different} architectures). We train the student models by utilizing an offline KD strategy \cite{gou2021knowledge}, where we distill knowledge from the more competent, pre-trained teacher model to the simpler, more lightweight student model. We augment the training of each student model with our KD protocols and evaluate the effect on predictive accuracy against the two vanilla KD approaches. 
If not mentioned otherwise, we utilize the following setup: we use MSE as a distillation loss $\loss_{\text{feat}}$; a learned linear layer as a transformation function $\mathcal{M}_s$ on the features of the student; and the identity transformation as $\mathcal{M}_t$. When distilling knowledge from GemNet-OC, we use the aggregated output block features $\boldsymbol{x}_\text{E}$ and $\boldsymbol{x}_\text{F}$ as node features and edge features, respectively. This is similar to the review-based setup proposed in \cite{chen2021distilling}. For PaiNN and SchNet, we use the final node features. Other training details are available in \cref{ap:train_hyperparams}.

% Main results table OC20
\begin{table}
\centering
\caption{Evaluation of the performance of our KD strategies across teacher-student architectures on the OC20 S2EF validation set. Best results are given in \textbf{bold}. All models have been trained on the OC20-2M dataset. Numbers in brackets represent the proportion of the gap between the student and the teacher that has been closed by the respective KD strategy (in \%). The validation results presented here are averaged across the 4 validation datasets available in OC20. Results for individual validation datasets are provided in \cref{ap:main_results_full}. }
\begin{tabular}{llcccccc}
\multicolumn{1}{c}{\multirow{3}{*}{}} & & & \multicolumn{5}{c}{OC20 S2EF Validation} 
\\\cmidrule(r){4-7}
\multicolumn{1}{c}{}   &   &                   &      Energy MAE  &    Force MAE  & Force cos & EFwT \\
&   \multicolumn{1}{l}{Model}     &        &      \si{\milli\electronvolt} $\downarrow$   &    \si{\milli\electronvolt\per\angstrom} $\downarrow$     & $\uparrow$  & \% $\uparrow$ \\\cmidrule(r){1-7}
% & \multicolumn{1}{l}{Model} & &   &   \multicolumn{2}{c}{OC20-2M} & & \multicolumn{2}{c}{Other dataset} &  \\\cmidrule(r){5-6}\cmidrule(r){8-9}
% \multicolumn{9}{c}{} \\
\parbox[t]{2mm}{\multirow{6}{*}{\rotatebox[origin=c]{90}{\emph{same}}}} & Student (PaiNN-small)               &   &              $489$      &     $47.1$     &  $0.345$ &  $0.085$  \\
& Teacher (PaiNN-big)              &  &          $440$        &   $45.3$     &   $0.376$      &   $0.139$  \\
& Vanilla KD (1)         &   &            $515 (\text{-}52.4\%)$          &          $48.5 (\text{-}81.0\%)$     &     $0.269 (\text{-}237\%)$        &  $0.07 (\text{-}28\%)$   \\
& Vanilla KD (2)         &   &                $476 (27.2\%)$      &          $50.8 (\text{-}215\%)$     &     $0.307 (\text{-}117\%)$        &   $0.068 (\text{-}32.6\%)$ \\
& n2n              &   &             $457 (\textbf{64.8\%})$         &          $46.7 (\textbf{20.5\%})$     &     $0.348 (\textbf{9.3\%})$        & $0.085 (\textbf{0.5\%})$    \\
& v2v                &   &          $459 (60.8\%)$            &          $47.2 (\text{-}9.1\%)$     &     $0.347 (6.8\%)$        & $0.079 (\text{-}11.9\%)$    \\
\multicolumn{7}{c}{} \\
\parbox[t]{2mm}{\multirow{5}{*}{\rotatebox[origin=c]{90}{\emph{similar}}}} & Student (SchNet)                  &   &              $1308$     &   $65.1$    &  $0.204$  &  $0$   \\
& Teacher (PaiNN-big)                       &   &              $440$        &   $45.3$     &   $0.376$      &   $0.139$    \\
& Vanilla KD (1)             &   &            $1214 (\textbf{10.8\%})$          &          $64.6 (2.3\%)$     &     $0.230 (\textbf{15.2\%})$        &  $0.003 (\textbf{1.8\%})$ \\
& Vanilla KD (2)               &   &            $1216 (10.5\%)$          &          $64.6 (\textbf{2.5\%})$     &     $0.229 (14.5\%)$        &  $0 (0\%)$ &  \\
& n2n              &   &            $1251 (6.6\%)$          &          $65.2 (\text{-}0.5\%)$     &     $0.223 (11.1\%)$        &  $0 (0\%)$  &  \\
\multicolumn{7}{c}{} \\
\parbox[t]{2mm}{\multirow{7}{*}{\rotatebox[origin=c]{90}{\emph{different}}}} & Student (PaiNN-big)                  &   &              $440$        &   $45.3$     &   $0.376$      &   $0.139$    \\
& Teacher (GemNet-OC)                      &   &              $286$        &   $25.7$     &   $0.598$      &   $1.063$     \\
& Vanilla KD (1)           &   &            $445 (\text{-}3.4\%)$          &          $44.3 (5.4\%)$     &     $0.372 (\text{-}2.2\%)$        &  $0.14 (0.1\%)$    \\
& Vanilla KD (2)             &   &            $419 (13.6\%)$          &          $114.8 (\text{-}353\%)$     &     $0.324 (\text{-}23.8\%)$        &  $0.127 (\text{-}1.3\%)$   \\
& n2n              &   &            $349 (\textbf{59.1\%})$          &          $43.5 (9.3\%)$     &     $0.379 (1.3\%)$        &  $0.227 (\textbf{9.5\%})$   \\
& e2n               &   &            $439 (0.6\%)$          &          $42.1 (\textbf{16.7\%})$     &     $0.391 (6.8\%)$        &  $0.186 (5.1\%)$     \\
& v2v              &   &            $439 (0.8\%)$          &          $42.3 (15.6\%)$     &     $0.394 (\textbf{7.9\%})$        &  $0.118 (\text{-}2.3\%)$        
% & e2e              &   &            $X (x\%)$          &          $X (x\%)$     &     $Y (y\%)$        &  $X (x\%)$  &          &       
\end{tabular}
\label{tab:main_results}
\end{table}


\textbf{Knowledge distillation results.} 
The results of the experiment are summarized in Table \ref{tab:main_results}, presenting a comparative analysis of the predictive performance of different student models with and without the implementation of knowledge distillation.
% A summary of the performance with and without knowledge distillation for the different teacher-student pairs can be found in Table \ref{tab:main_results}. The results indicate that our feature-based KD strategies are effective across all teacher-student categories. 
% 
Especially notable was the improvement in energy predictions for PaiNN-big when trained with GemNet-OC as a teacher under the \emph{n2n} KD strategy. This configuration led to a reduction in PaiNN-big's energy MAE from 440 \si{\milli\electronvolt} to 349 \si{\milli\electronvolt}. This equates to a decrease of over 20 \% compared to the base PaiNN-big model, thus closing nearly 60 \% of the accuracy gap between GemNet-OC and PaiNN-big in energy predictions. The \emph{n2n} KD strategy also led to an improvement in force predictions, albeit a less dramatic one, bridging 9.3 \% of the performance gap between the two models. We manage to improve that with our custom \emph{e2n} and \emph{v2v} protocols, elevating the improvements in force predictions to 16.7 \% and 15.6\% respectively. Nonetheless, this was offset by a decline in energy prediction accuracy, which dropped to 0.6 \% and 0.8 \% respectively. 
% 
In contrast, the \emph{n2n} proved to be the best strategy for both energy and force predictions when distilling PaiNN-big into PaiNN-small. For this teacher-student configuration, we were able to close the gap in energy and force prediction between the two models by ${\sim}65\%$ and $20\%$ respectively. The effect of \emph{v2v} KD on energy predictions was similar at 60.8 \%. Yet, this approach was detrimental to force predictions (-9.8 \%).
% 
In both cases, the performance of our two vanilla KD protocols was not as effective, either yielding smaller improvements or significantly impairing accuracy. Still, they showed the most promising results when distilling knowledge from PaiNN-big to SchNet. However, it is noteworthy to highlight that the effect of KD for this pair was tangibly lower at around 11 \% for energy improvements and 2.5 \% for forces. Based on previous research \cite{cho2019efficacy}, we deduct this may be a result of the vast gap in complexity and performance between the two models, which affects the learning of the substantially more lightweight student model.
 
These experiments showcase the effectiveness of our KD strategies, which allowed us to consistently boost the accuracy of student models across all examined teacher-student configurations. Importantly, these findings remained consistent when broadening our experiment to the COLL dataset, where we similarly observed significant performance gains across all model configurations (see \cref{ap:coll}). This further highlights the robustness of our training protocols and their utility across molecular datasets. Finally, we reiterate a central facet of our approach in that we enhance the performance of our student models solely by imposing an additional KD loss during training, without any structural alterations to their architecture. Hence, % despite increasing training times (we need to perform additional forward passes through the teacher during training), 
the inference throughput of the distilled student models remained completely unaffected.

\textbf{Effect of KD on model similarity.}
We continue our investigation of KD on molecular GNNs by investigating how KD affects the student architectures in the various teacher-student configurations we have studied. To this end, we analyze how the CKA similarity between the teacher and the student models changes with the introduction of KD during training. We observed that KD introduces strong and specific similarity gains in layers that are used for KD, which also propagate along the student architecture (see \cref{fig:cka_after_kd}). We saw similar behavior across KD strategies and teacher-student configurations
(see \cref{ap:cka}), allowing us to monitor and quantify the effect of KD on student models as we explore different KD settings and design choices.
% This allowed us to monitor the effect of KD on the student model as we  KD strategies and teacher-student configurations (see \cref{ap:cka}).
% , as well as trends demonstrating that higher similarity gains are reflective of higher boosts in performance. This demonstrates the effectiveness of the approach as a tool for profiling KD procedures.
Hence, we believe that such similarity-based analyses can serve as a useful tool for quantitatively profiling KD procedures, which can elucidate the mechanism behind KD and inform the development of KD protocols. 

\begin{figure}[h]
  \centering
\includegraphics[width=0.9\textwidth]{cka/after_n2n.png}
  \caption{Similarity analysis between Gemnet-OC and PaiNN-big without KD (left) and with KD (right). It illustrates a clear similarity gain within the feature pair used during KD (indicated with a $\textcolor{cyan_our}{\swarrow}$ in the middle heatmap), which also propagates to earlier layers of the student architecture. Similarity analyses for other KD strategies and teacher-student configurations are presented in \cref{ap:cka}.} 
  \label{fig:cka_after_kd}
\end{figure} 

\textbf{Hyperparameter studies} 
To increase the breadth of our work, we additionally conducted a thorough study to evaluate the effect of different design choices within our KD strategy. We present the result of this study here, hoping to provide valuable insights into the trends we have observed, which can inform future research in the area. The results are further detailed in \cref{ap:hyperparams}.

\emph{Effect of distillation loss}: In terms of distillation losses $\loss_{\text{feat}}$, we explored: mean squared error (MSE); mean absolute error (MAE); and also directly optimizing CKA. Our analyses showed that MAE performed similarly to MSE, whereas using CKA substantially hurt accuracy.

\emph{Effect of transformation function}: 
We also investigated a number of different transformation functions and the effect they have on performance. The transformations we utilized include: the identity transformation (when appropriate); learned linear transformations and MLP projection heads; as well as variants of the structure-preserving transformations described in \cite{joshi2022}. Our results (see \cref{ap:transformations}) showed that the examined transformations either hurt performance or only provided a marginal improvement in performance compared to our default linear mapping.

\emph{Effect of feature selection}: We additionally explored the effect of feature selection on KD performance. In particular, we analyzed the change in the predictive accuracy of PaiNN as a student model as we distill: features from earlier layers in the teacher (GemNet-OC); onto earlier layers in the student; across multiple layers in the teacher and the student. We observed that our default, review-based strategy was the best-performing KD protocol for this teacher-student configuration, with other approaches either hurting performance or only providing a marginal improvement (see \cref{ap:features}). We complemented our experiments by performing CKA-based similarity analyses, which allowed us to gain insights into how these protocols change the architecture of the student (see \cref{ap:cka}).

%\emph{Effect of strategy combinations}: Lastly, we study the effect of combining losses. For instance, we compare best n2n to best n2n + v2v ... Results in \cref{ap:combinations} ... we see that ... CKA analysis allowed us to monitor changes in the similarities of the models ...

\section{Data augmentation}
As for most other applications, data labeling for molecular data is costly as it requires running computationally expensive quantum mechanical calculations to obtain energies and forces. 
%From a knowledge distillation perspective, the lack of labeled data could be overcome using the teacher to label unlabelled data, thus distilling its knowledge via additional data. Additionally, as atoms can move in a continuum and atom types are replaceable, it is relatively easy to generate new molecular systems. 
We explored two protocols for knowledge distillation via training on novel data points labeled by the teacher, and briefly describe them here.
See \cref{ap:data_aug} for more details. 

% These include data perturbation, jittering, synthetic data, sample weighing/masking, pertaining?

\textbf{Random rattling.}
Apart from the relaxation trajectories in OC20, there is also a dataset called "Rattled", where random noise has been added to a subset of the systems along the relaxation trajectories, and the forces and energies have then been computed for these rattled systems using DFT. Adding noise to existing structures has also been used in the context of pretraining of molecular GNNs \cite{zaidi2022pre, magar_crystal_2022} or as a regularization strategy \cite{godwin2021simple}. 
%Inspired by this, we added random noise to the atomic positions and then used the teacher to provide energy and force labels for these rattled samples. 
However, adding noise to existing samples and aligning the predictions between teacher and student (without ground truth labels) did not show significant improvement.
%This approach showed mixed results. For example, pairing it with the v2v loss improved the energy predictions by $9\%$ compared to using regular \emph{v2v}, but made the force predictions worse, while pairing it with n2n worsened both metrics. 
Additionally, we tried using gradient ascent to find perturbations that maximize the discrepancy between the teacher and student predictions, similar to \cite{schwalbe-koda_differentiable_2021}, but this approach increased training time and did not show improvements over random noise. 

\textbf{Synthetic Data.}
%Samples in the OC20 datasets are generated by first running DFT-based relaxations on randomly sampled systems, i.e., a combination of adsorbate and catalyst, and subsequently randomly sampling individual frames along these relaxation trajectories. 
One weakness of the OC20 datasets is that some samples originate from the same relaxation trajectory and are therefore correlated. %By sampling along the same trajectory, many samples are correlated, which limits the diversity of the dataset. 
To tackle this, % issue and explore the capabilities of knowledge distillation, 
we generated our own distilled dataset coined \textit{d1M} consisting of 1M samples, which were generated by sampling novel systems (generated with the OC Datasets codebase \footnote{\url{https://github.com/Open-Catalyst-Project/Open-Catalyst-Dataset}}), running relaxations with the pre-trained GemNet-OC 2M model, and then subsampling approx. 10\% of the frames.
We tried different ways of incorporating this new d1M dataset, all of which are based on joint training on the 2M as well as the d1M datasets (similar to what we did with the rattled systems). 
%To explore different dataset combinations, we introduce two new hyperparameters: the ratio of samples drawn from the d1M dataset $\alpha_{d1M}$, and the loss weighting ratio between DFT- and GemNet-OC-based samples $r_{d1M/DFT}$. However, after extensively exploring these two hyperparameters for training our PaiNN-big model on the 2M+d1M data, we could not see improvement over the baseline performance. This result is surprising as similar approaches have been used for knowledge distillation in other domains \cite{zhang2020pushing}. We hypothesize that by the out-of-distribution nature of the systems in the d1M dataset, some of the labels from the pre-trained GemNet-OC model might lead the student's training in the wrong direction. For completeness, we also explored knowledge distillation on a combination of both datasets on the n2n as well as v2v distillation from GemNet-OC to PaiNN-big, but did not observe significant improvement either.
To explore different combinations of the ground truth DFT samples and the d1M samples during training, we defined two hyperparameters determining (a) how many of the samples per batch should originate from each of the datasets, and (b) how to weight the loss contributions based on the dataset of origin. Unfortunately, and contrary to similar approaches e.g. in speech recognition \cite{zhang2020pushing}, the results we observed didn't significantly improve on the baseline models. 

% We generated 1M synthetic samples by first drawing 100k random adsorbate and catalyst combinations (systems) and then running relaxations with a pre-trained GemNet-OC model. Out of these relaxations with 100 steps on average (200 max), we randomly draw approx. 10\% to obtain 1M samples.

% In the next step, we combine the 1M samples with the 2M OCP dataset, which is based on DFT relaxations. Directly working with this combined dataset means iterating over a 1-to-2 ratio of samples from each subset in an epoch of 3M samples. To control this ratio, we define the target ratio of samples from the synthetic dataset during training $\alpha_{target}\in[0,1]$. Setting  $\alpha_{target}=0.5$ means that per epoch we iterate over the 1M dataset 1.5 times and over the 2M dataset 0.75 times.


% Samples in the OC20 datasets are generated by first running DFT-based relaxations on randomly sampled systems, i.e. combination of adsorbate and catalyst, and subsequently randomly sampling individual frames along these relaxation trajectories. By sampling along the same trajectory, many samples are correlated, which limits the diversity of the dataset. To tackle this issue and explore the capabilities of knowledge distillation, we generated our own distilled 1M (d1M) dataset by sampling novel systems (generated with the OC Datasets codebase \footnote{\url{https://github.com/Open-Catalyst-Project/Open-Catalyst-Dataset}}) and running relaxations with a pre-trained GemNet-OC 2M model \footnote{\url{https://github.com/Open-Catalyst-Project/ocp/blob/main/MODELS.md}}.

% We tried different protocols for incorporating this new d1M dataset, all of which are based on jointly training on the 2M as well as the d1M datasets. To explore different dataset combinations, we introduce two new hyperparameters: the ratio of samples drawn from the d1M dataset $\alpha_{d1M}$, and the loss weighting ratio between DFT- and GemNet-OC-based samples $r_{d1M/DFT}$. However, after extensively exploring these two hyperparameters for training our PaiNN-big model on the 2M+d1M data, we could not see improvement over the baseline performance. This result is surprising as similar approaches have been used for knowledge distillation in other domains \cite{xu2022contrastive}, and we hypothesize that by the out-of-distribution nature of the systems in the d1M dataset, some of the labels from the pre-trained GemNet-OC model might lead the student's training in the wrong direction. For completeness, we also explored knowledge distillation on a combination of both datasets on the n2n as well as v2v distillation from GemNet-OC to PaiNN-big, but didn't observe significant improvement either.


\section{Conclusion}
In this paper, we investigate the utility of knowledge distillation as a means of distilling larger, more computationally expensive GNNs for molecules into smaller, more efficient models. To this end, we propose three distinct feature-based KD strategies, which facilitate the distillation of intermediate representations across diverse GNN models. Our analyses demonstrate that our protocols can significantly enhance the performance of molecular GNN without any modifications to their architecture, allowing us to run faster molecular simulations without substantially impairing predictive accuracy. We validate our approach across different teacher-student configurations and datasets, showcasing the effectiveness and robustness of the method. One limitation of our approach is that even though inference times are not affected, we need to perform forward passes through the teacher during training, increasing training times.

With this work, we aim to elucidate the potential of KD in the domain of molecular simulations and stimulate future work in the area, which can bring about important advances in scientific disciplines as broad as material science, drug discovery and catalysis. However, it is noteworthy to highlight that such technologies can be potentially harmful if, e.g., used to simulate and discover systems that are toxic or intended to be used in harmful technology.

% confirming that KD is a useful technique for addressing the pressing trade-off between predictive accuracy and computational complexity of modern GNNs for molecules.
% With this work, we want to highlight the utility of knowledge distillation in enhancing the performance of molecular GNNs and accelerating molecular simulations. 

% The use of molecular GNNs can help to speed up molecular simulations, which find applications in important scientific disciplines, such as material science, drug discovery and catalysis. 

% To fully benefit from using GNNs for molecular simulations, there is a need to keep the computational time to a minimum. Therefore, we have utilised knowledge distillation to distill larger, more computationally expensive GNNs for molecules into smaller and computationally faster models. 
% We show how the performance of the smaller models can be improved without any modifications of their architecture, meaning that their inference times are unaffected. In many cases, we substantially close the gap in performance between teacher and student. One limitation of this approach is that even though the inference times are not affected, during training, we need to perform forward passes through the teacher. With this work, we want to highlight the possibility of using knowledge distillation as a mean of improving molecular GNNs. We believe that a lot more can be done, and we hope this paper will serve as a foundation and motivation for future investigation into the area. The use of molecular GNNs can help to speed up molecular simulations, which finds applications in important subjects like material science, drug discovery, etc. However, these applications can also be potentially harmful if the models are used to simulate and discover systems that are, e.g., toxic.

\begin{ack} % Acknowledgements and disclosure of funding
% computing resources

This project would not have been possible without the computing resources provided by: the Knut and Alice Wallenberg Foundation at the National Supercomputer Centre (Berzelius resource); the National Academic Infrastructure for Supercomputing in Sweden (NAISS) at Chalmers e-Commons at Chalmers (C3SE) partially funded by the Swedish Research Council through grant agreement no. 2022-06725; as well as the Chair of Aerodynamics and Fluid Mechanics at Technical University of Munich. 
% The computations have been enabled by Berzelius resource provided by the Knut and Alice Wallenberg Foundation at the National Supercomputer Centre, resources provided by the National Academic Infrastructure for Supercomputing in Sweden (NAISS) at Chalmers e-Commons at Chalmers (C3SE) partially funded by the Swedish Research Council through grant agreement no. 2022-06725, as well as the Chair of Aerodynamics and Fluid Mechanics at Technical University of Munich. 
% LOGML and previous members: 
We are also grateful to the team behind 2022 London Geometry and Machine Learning Summer School (LOGML), where this research project was initially conceived. We would like to specially thank Guocheng Qian and I-Ju Chen for their contribution during the early conceptualizing stages of this project during and in the first weeks following the summer school. 
% OC team
We also thank the Open Catalyst team for their open-source codebase, support and discussions. In particular, we would like to thank Muhammed Shuaibi for providing the COLL dataset in LMDB format.
% Other
Figures assembled in BioRender.com.


% funding: 
Funding disclosure: 
Dimitar Georgiev is supported by UK Research and Innovation [UKRI Centre for Doctoral Training in AI for Healthcare grant number EP/S023283/1];
% 
Filip Ekström Kelvinius is financially supported by the Excellence Center at Linköping--Lund in Information Technology (ELLIIT).
% 

\end{ack}


\bibliography{main}
\bibliographystyle{unsrt}

\newpage
\appendix

\section{Training and hyperparameters} 
\label{ap:train_hyperparams}
The models were trained on 40 GB A100s and 32 GB RTX600 GPUs. All models were trained on single GPUs, except for SchNet when trained on OC20, which required 3 GPUs.
In Tables~\ref{tab:GemNet_hypers}-\ref{tab:schnet_hypers}, hyperparameters for the models are presented. In Table~\ref{tab:KD_hypers}, the weighting of the KD are presented.
\begin{table}[H]
\centering
\caption{GemNet-OC hyperparameters. For the experiments we have conducted on OC20 data, we used the model pretrained on OC20-2M (July 2022), and we have hence not done any training ourselves.}
\begin{tabular}{lcc}
\label{tab:GemNet_hypers}
Hyperparameter & OC20 & COLL\\\hline
No. spherical basis & 7 & 7 \\
No. radial basis & 128 & 128\\
No. blocks & 4 & 4\\
Atom embedding size & 256 & 128 \\
Edge embedding size & 512 & 256 \\\\

Triplet edge embedding input size & 64 & 64\\
Triplet edge embedding output size & 64 & 64\\
Quadruplet edge embedding input size & 32 & 32\\
Quadruplet edge embedding output size & 32 & 32\\
Atom interaction embedding input size & 64 & 64\\
Atom interaction embedding output size & 64 & 64\\
Radial basis embedding size & 16 & 16\\
Circular basis embedding size & 16 & 16  \\
Spherical basis embedding size & 32 & 32 \\\\

No. residual blocks before skip connection & 2 & 2\\
No. residual blocks after skip connection & 2 & 2\\
No. residual blocks after concatenation & 1 & 1 \\
No. residual blocks in atom embedding blocks & 3 & 3 \\
No. atom embedding output layers & 3 & 3 \\\\

Cutoff & 12.0 & 12.0\\
Quadruplet cutoff & 12.0 &  12.0\\
Atom edge interaction cutoff &  12.0 & 12.0 \\
Atom interaction cutoff & 12.0 & 12.0 \\
Max interaction neighbors & 30 & 30\\
Max quadruplet interaction neighbors & 8 & 8 \\
Max atom edge interaction neighbors & 20 & 20\\
Max atom interaction neighbors & 1000 & 1000\\\\

Radial basis function & Gaussian & Gaussian\\
Circular basis function & Spherical harmonics & Spherical Harmonics\\
Spherical basis function & Legendre Outer & Legendre Outer\\
Quadruplet interaction & True & True\\
Atom edge interaction & True & True\\
Edge atom interaction & True & True\\
Atom interaction & True & True \\
Direct forces & True & True \\\\

Activation & Silu & Silu \\
Optimizer & - & AdamW\\
Scheduler & - & LinearWarmupExponentialDecay \\
Force coefficient & - & 100 \\
Energy coefficient & - & 1 \\
EMA decay & - & 0.999\\
Gradient clip norm threshold & - & 10\\
Initial learning rate & - & $10^{-3}$\\
\end{tabular}
\end{table}

\begin{table}[H]
\centering
\caption{PaiNN hyperparameters. The different number of layers and hidden channels indicate the values we have used for PaiNN-big and PaiNN-small, respectively.}
\begin{tabular}{lcc}
Hyperparameter & OC20 & COLL\\\hline
Hidden channels & 512/256 & 256/128 \\
Number of layers & 6/4 & 6/4\\
Number of RBFs & 128 & 128\\
Cutoff & 12.0 & 12.0\\
Max. num. neighbors & 50 & 50\\
Direct Forces & True & True\\\\

Batch size & 32 & 32\\
Optimizer & AdamW & AdamW\\
AMSGrad & True & True\\
Initial learning rate & $10^{-4}$ & $10^{-3}$\\
Scheduler & LambdaLR & LinearWarmupExponentialDecay\\
Warmup steps & None & 3750\\
Learning rate decay factor & 0.45 & 0.01 \vspace{0.03in}\\
Learning rate milestones (steps) & \begin{tabular}[c]{@{}c@{}}160000, 320000, \\ 480000, 640000\end{tabular} & - \vspace{0.03in}\\
Force coefficient & 100 & 100 \\
Energy coefficient & 1 & 1\\
EMA decay & 0.999 & 0.999\\
Gradient clip norm threshold & 10 & 10\\
\end{tabular}
\label{tab:painn_hypers}
\end{table}

\begin{table}[H]
\centering
\caption{SchNet hyperparameters. For the experiments we have conducted on OC20 data, we used the default configuration for the OCP-2M dataset as provided in the OCP repository.}
\begin{tabular}{lcc}
\label{tab:schnet_hypers}
Hyperparameter & OC20 & COLL\\\hline
Hidden channels & 1024 & 128 \\
Filters & 256 & 128 \\
Interaction blocks & 5 & 6\\
Gaussians & 200 & 50\\
Cutoff & 6.0 & 12.0\\\\
Batch size & 192 & 32\\
Initial learning rate & $10^{-4}$ & $10^{-3}$\\
Optimizer & AdamW &  AdamW \\
Scheduler & LambdaLR & LinearWarmupExponentialDecay \\

Learning rate decay factor & 0.1 & 0.01\vspace{0.03in}\\
Learning rate milestones & \begin{tabular}[c]{@{}c@{}}52083, 83333, \\ 104166\end{tabular} & - \vspace{0.03in} \\
Warmup steps & 31250 & 3750\\
Warmup factor & 0.1 & - \\
Force Coefficient & 100 & 100 \\
Energy Coefficient & 1 &  1\\
\end{tabular}
\end{table}

\begin{table}[H]
\centering
\caption{Choice of the weighting factor $\lambda$ of the KD loss for the different teacher-student configurations and KD strategies.}
\begin{tabular}{lcccc}
\label{tab:KD_hypers}
Teacher & Student & Loss & OC20 & COLL\\\hline
GemNet-OC & PaiNN-big & vanilla (1) & $1.0$ & $0.2$\\
GemNet-OC & PaiNN-big & vanilla (2) & $500$ & $100$ \\
GemNet-OC & PaiNN-big & n2n         & $10000$ & $1000$ \\
GemNet-OC & PaiNN-big & e2n         & $10$ & $10$ \\
GemNet-OC & PaiNN-big & v2v         & $50000$ & $100$ \\\\

PaiNN-big & PaiNN-small & vanilla (1) & $1$ & $1$ \\
PaiNN-big & PaiNN-small & vanilla (2) & $200$ & $100$ \\
PaiNN-big & PaiNN-small & n2n & $100$ & $100$ \\
PaiNN-big & PaiNN-small & v2v & $1000$& $10000$ \\\\

PaiNN-big & SchNet & vanilla (1) & $0.1$ & $1$ \\
PaiNN-big & SchNet & vanilla (2) & $0.1$ & $100$ \\
PaiNN-big & SchNet & n2n & $1000$ & $100$ \\
\end{tabular}
\end{table}

\section{Full validation results} 
\label{ap:main_results_full}

Tables \ref{tab:val-id}-\ref{tab:ood_both_results} show the results achieved with the baseline KD setup (\cref{tab:main_results}) on the 4 separate validation sets in the OC20 S2EF.

\begin{table}[H]
\centering
\caption{Evaluation results on the OC20 S2EF \textbf{in-distribution} validation set.}
\begin{tabular}{llcccccc}
\multicolumn{1}{c}{\multirow{3}{*}{}} & & & \multicolumn{5}{c}{OC20 S2EF Validation (\textbf{in-distribution})} 
\\\cmidrule(r){4-7}
\multicolumn{1}{c}{}   &   &                   &      Energy MAE  &    Force MAE  & Force cos & EFwT \\
&   \multicolumn{1}{l}{Model}     &        &      \si{\milli\electronvolt} $\downarrow$   &    \si{\milli\electronvolt\per\angstrom} $\downarrow$     & $\uparrow$  & \% $\uparrow$ \\\cmidrule(r){1-7}
% & \multicolumn{1}{l}{Model} & &   &   \multicolumn{2}{c}{OC20-2M} & & \multicolumn{2}{c}{Other dataset} &  \\\cmidrule(r){5-6}\cmidrule(r){8-9}
% \multicolumn{9}{c}{} \\
\parbox[t]{2mm}{\multirow{6}{*}{\rotatebox[origin=c]{90}{\emph{same}}}} & Student (PaiNN-small)               &   &              $409$      &     $41.6$     &  $0.357$ &  $0.16$  \\
& Teacher (PaiNN-big)              &  &          $358$        &   $38.5$     &   $0.390$      &   $0.25$  \\
& Vanilla KD (1)         &   &            $426 (\text{-}33.7\%)$          &          $42.6 (\text{-}32.3\%)$     &     $0.35 (\text{-}21.9\%)$        &  $0.12 (\text{-}44.4\%)$   \\
& Vanilla KD (2)         &   &                $396 (25.7\%)$      &          $45.7 (\text{-}132.3\%)$     &     $0.316 (\text{-}126.9\%)$        &   $0.11 (\text{-}55.6\%)$ \\
& n2n              &   &             $393 (\textbf{31.2\%})$         &          $41.7 (\text{-}2.5\%)$     &     $0.359 (\textbf{4.7\%})$        & $0.15 (\text{-}7.8\%)$    \\
& v2v                &   &          $406 (5.6\%)$            &          $42.1 (\text{-}16.9\%)$     &     $0.357 (\text{-}2.6\%)$        & $0.13 (\text{-}32.6\%)$    \\
\multicolumn{7}{c}{} \\
\parbox[t]{2mm}{\multirow{5}{*}{\rotatebox[origin=c]{90}{\emph{similar}}}} & Student (SchNet)                  &   &              $1237$     &   $62.2$    &  $0.214$  &  $0$   \\
& Teacher (PaiNN-big)              &  &          $358$        &   $38.5$     &   $0.390$      &   $0.25$  \\
& Vanilla KD (1)             &   &            $1139 (\textbf{10.6\%})$          &          $52.9 (\textbf{13.1\%})$     &     $0.2422 (\textbf{16\%})$        &  $0 (0\%)$ \\
& Vanilla KD (2)               &   &            $1140 (10.5\%)$          &          $59.2 (12.7\%)$     &     $0.241 (15.1\%)$        &  $0 (0\%)$ &  \\
& n2n              &   &            $1170 (7\%)$          &          $60 (9.3\%)$     &     $0.235 (11.9\%)$        &  $0 (0\%)$  &  \\
\multicolumn{7}{c}{} \\
\parbox[t]{2mm}{\multirow{7}{*}{\rotatebox[origin=c]{90}{\emph{different}}}} & Student (PaiNN-big)              &  &          $358$        &   $38.5$     &   $0.390$      &   $0.25$  \\
& Teacher (GemNet-OC)                      &   &              $226$        &   $22.5$     &   $0.61$      &   $1.89$     \\
& Vanilla KD (1)           &   &            $364 (\text{-}3.9\%)$          &          $38.9 (\text{-}2.4\%)$     &     $0.385 (\text{-}2.1\%)$        &  $0.25 (0.1\%)$    \\
& Vanilla KD (2)             &   &            $357 (0.7\%)$          &          $43.5 (\text{-}31.4\%)$     &     $0.334 (\text{-}25.4\%)$        &  $0.21 (\text{-}2.4\%)$   \\
& n2n              &   &            $276 (\textbf{62.2\%})$          &          $38.3 (1.4\%)$     &     $0.385 (\text{-}2\%)$        &  $0.405 (\textbf{9.5\%})$   \\
& e2n               &   &            $341 (13.1\%)$          &          $37.7 (\textbf{8.6\%})$     &     $0.405 (6.9\%)$        &  $0.316 (4\%)$     \\
& v2v              &   &            $373 (\text{-}11.3\%)$          &          $37.5 (6.1\%)$     &     $0.406 (\textbf{7.3\%})$        &  $0.209 (\text{-}2.5\%)$        
% & e2e              &   &            $X (x\%)$          &          $X (x\%)$     &     $Y (y\%)$        &  $X (x\%)$  &          &       
\end{tabular}
\label{tab:val-id}
\end{table}

\begin{table}[H]
\centering
\caption{Evaluation results on the OC20 S2EF \textbf{out-of-distribution (adsorbates)} validation set.}
\begin{tabular}{llcccccc}
\multicolumn{1}{c}{\multirow{3}{*}{}} & & & \multicolumn{5}{c}{OC20 S2EF Validation (\textbf{out-of-distribution (adsorbates)})} 
\\\cmidrule(r){4-7}
\multicolumn{1}{c}{}   &   &                   &      Energy MAE  &    Force MAE  & Force cos & EFwT \\
&   \multicolumn{1}{l}{Model}     &        &      \si{\milli\electronvolt} $\downarrow$   &    \si{\milli\electronvolt\per\angstrom} $\downarrow$     & $\uparrow$  & \% $\uparrow$ \\\cmidrule(r){1-7}
% & \multicolumn{1}{l}{Model} & &   &   \multicolumn{2}{c}{OC20-2M} & & \multicolumn{2}{c}{Other dataset} &  \\\cmidrule(r){5-6}\cmidrule(r){8-9}
% \multicolumn{9}{c}{} \\
\parbox[t]{2mm}{\multirow{6}{*}{\rotatebox[origin=c]{90}{\emph{same}}}} & Student (PaiNN-small)               &   &              $469$      &     $47.9$     &  $0.334$ &  $0.03$  \\
& Teacher (PaiNN-big)                  &   &              $437$        &   $44.5$     &   $0.369$      &   $0.043$    \\
& Vanilla KD (1)         &   &            $519 (\text{-}156.2\%)$          &          $47.4 (14.7\%)$     &     $0.334 (\text{-}1.2\%)$        &  $0.003 (0\%)$   \\
& Vanilla KD (2)         &   &                $477 (\text{-}23.6\%)$      &          $54.5 (\text{-}106.2\%)$     &     $0.297 (\text{-}109.9\%)$        &   $0.002 (\text{-}76.9\%)$ \\
& n2n              &   &             $443 (80.9\%)$         &          $47.2 (\textbf{19.8\%})$     &     $0.337 (\textbf{9.3\%})$        & $0.003 (\text{-}10\%)$    \\
& v2v                &   &          $441 (\textbf{87.5}\%)$            &          $47.9 (\text{-}1\%)$     &     $0.337 (8.68\%)$        & $0.04 (\textbf{81.6\%})$    \\
\multicolumn{7}{c}{} \\
\parbox[t]{2mm}{\multirow{5}{*}{\rotatebox[origin=c]{90}{\emph{similar}}}} & Student (SchNet)                  &   &              $1344$     &   $58$    &  $0.196$  &  $0$   \\
& Teacher (PaiNN-big)                  &   &              $437$        &   $44.5$     &   $0.369$      &   $0.043$    \\
& Vanilla KD (1)             &   &            $1247 (10.7\%)$          &          $64.5 (\text{-}49.7\%)$     &     $0.221 (\textbf{14.7\%})$        &  $0 (0\%)$ \\
& Vanilla KD (2)               &   &            $1245 (\textbf{10.9\%})$          &          $64.5 (\text{-}48.2)$     &     $0.22 (14\%)$        &  $0 (0\%)$ &  \\
& n2n              &   &            $1286 (6.3\%)$          &          $65 (\text{-}51.9\%)$     &     $0.213 (9.9\%)$        &  $0 (0\%)$  &  \\
\multicolumn{7}{c}{} \\
\parbox[t]{2mm}{\multirow{7}{*}{\rotatebox[origin=c]{90}{\emph{different}}}} & Student (PaiNN-big)                  &   &              $437$        &   $44.5$     &   $0.369$      &   $0.043$    \\
& Teacher (GemNet-OC)                      &   &              $258$        &   $25.2$     &   $0.6$      &   $0.45$     \\
& Vanilla KD (1)           &   &            $438 (\text{-}0.8\%)$          &          $45.1 (\text{-}3.1\%)$     &     $0.362 (\text{-}2.6\%)$        &  $0.056 (3.1\%)$    \\
& Vanilla KD (2)             &   &            $408 (15.9\%)$          &          $49.2 (\text{-}24.3\%)$     &     $0.315 (\text{-}23.1\%)$        &  $0.036 (\text{-}1.7\%)$   \\
& n2n              &   &            $327 (\textbf{61.2\%})$          &          $44.1 (2.2\%)$     &     $0.374 (2.4\%)$        &  $0.084 (\textbf{10.1\%})$   \\
& e2n               &   &            $453 (\text{-}9.1\%)$          &          $42.4 (10.9\%)$     &     $0.384 (6.7\%)$        &  $0.063 (4.9\%)$     \\
& v2v              &   &            $421 (8.9\%)$          &          $424 (\textbf{11\%})$     &     $0.388 (\textbf{8.4\%})$        &  $0.054 (2.7\%)$        
% & e2e              &   &            $X (x\%)$          &          $X (x\%)$     &     $Y (y\%)$        &  $X (x\%)$  &          &       
\end{tabular}
\label{tab:ood_ads_results}
\end{table}

\begin{table}[H]
\centering
\caption{Evaluation results on the OC20 S2EF \textbf{out-of-distribution (catalysts)} validation set.}
\begin{tabular}{llcccccc}
\multicolumn{1}{c}{\multirow{3}{*}{}} & & & \multicolumn{5}{c}{OC20 S2EF Validation (\textbf{out-of-distribution (catalysts)})} 
\\\cmidrule(r){4-7}
\multicolumn{1}{c}{}   &   &                   &      Energy MAE  &    Force MAE  & Force cos & EFwT \\
&   \multicolumn{1}{l}{Model}     &        &      \si{\milli\electronvolt} $\downarrow$   &    \si{\milli\electronvolt\per\angstrom} $\downarrow$     & $\uparrow$  & \% $\uparrow$ \\\cmidrule(r){1-7}
% & \multicolumn{1}{l}{Model} & &   &   \multicolumn{2}{c}{OC20-2M} & & \multicolumn{2}{c}{Other dataset} &  \\\cmidrule(r){5-6}\cmidrule(r){8-9}
% \multicolumn{9}{c}{} \\
\parbox[t]{2mm}{\multirow{6}{*}{\rotatebox[origin=c]{90}{\emph{same}}}} & Student (PaiNN-small)               &   &              $467$      &     $42$     &  $0.341$ &  $0.13$  \\
& Teacher (PaiNN-big)                  &   &              $412$        &   $39.2$     &   $0.369$      &   $0.23$    \\
& Vanilla KD (1)         &   &            $466 (4.9\%)$          &          $42.8 (\text{-}28.4\%)$     &     $0.336 (\text{-}16.9\%)$        &  $0.11 (\text{-}20\%)$   \\
& Vanilla KD (2)         &   &                $439 (52.4\%)$      &          $45.4 (\text{-}120.6\%)$     &     $0.306 (\text{-}120.8\%)$        &   $0.12 (\text{-}10\%)$ \\
& n2n              &   &             $437 (\textbf{56.3\%})$         &          $42.0 (\textbf{1\%})$     &     $0.343 (\textbf{8.2\%})$        & $0.14 (\textbf{11.2\%})$    \\
& v2v                &   &          $444 (43.4\%)$            &          $42.4 (\text{-}12.8\%)$     &     $0.342 (3.6\%)$        & $0.12 (\text{-}12\%)$    \\
\multicolumn{7}{c}{} \\
\parbox[t]{2mm}{\multirow{5}{*}{\rotatebox[origin=c]{90}{\emph{similar}}}} & Student (SchNet)                  &   &              $1205$     &   $61.6$    &  $0.205$  &  $0$   \\
& Teacher (PaiNN-big)                  &   &              $412$        &   $39.2$     &   $0.369$      &   $0.23$    \\
& Vanilla KD (1)             &   &            $1122 (\textbf{10.6\%})$          &          $58.7 (\textbf{12.9\%})$     &     $0.234 (\textbf{15.9\%})$        &  $0.01 (\textbf{4.3\%})$ \\
& Vanilla KD (2)               &   &            $1122 (10.5\%)$          &          $58.8 (12.5\%)$     &     $0.23 (15.2\%)$        &  $0 (0\%)$ &  \\
& n2n              &   &            $1150 (6.9\%)$          &          $59.4 (9.8\%)$     &     $0.225 (12.3\%)$        &  $0 (0\%)$  &  \\
\multicolumn{7}{c}{} \\
\parbox[t]{2mm}{\multirow{7}{*}{\rotatebox[origin=c]{90}{\emph{different}}}} & Student (PaiNN-big)                  &   &              $412$        &   $39.2$     &   $0.369$      &   $0.23$    \\
& Teacher (GemNet-OC)                      &   &              $288$        &   $24$     &   $0.576$      &   $1.68$     \\
& Vanilla KD (1)           &   &            $420 (\text{-}6.5\%)$          &          $39.5 (\text{-}2.3\%)$     &     $0.365 (\text{-}1.8\%)$        &  $0.21 (\text{-}1.1\%)$    \\
& Vanilla KD (2)             &   &            $400 (9.5\%)$          &          $43.6 (\text{-}29\%)$     &     $0.320 (\text{-}23.5\%)$        &  $0.23 (0.2\%)$   \\
& n2n              &   &            $343 (\textbf{55.5\%})$          &          $39 (1.2\%)$     &     $0.374 (2.2\%)$        &  $0.36 (\textbf{9\%})$   \\
& e2n               &   &            $398 (11.1\%)$          &          $37.8 (\textbf{8.8\%})$     &     $0.383 (6.8\%)$        &  $0.32 (6.2\%)$     \\
& v2v              &   &            $425 (\text{-}10.7\%)$          &          $38.4 (5\%)$     &     $0.383 (\textbf{6.9\%})$        &  $0.17 (\text{-}4\%)$        
% & e2e              &   &            $X (x\%)$          &          $X (x\%)$     &     $Y (y\%)$        &  $X (x\%)$  &          &       
\end{tabular}
\label{tab:ood_cat_results}
\end{table}

\begin{table}[H]
\centering
\caption{Evaluation results on the OC20 S2EF \textbf{out-of-distribution (both)} validation set.}
\begin{tabular}{llcccccc}
\multicolumn{1}{c}{\multirow{3}{*}{}} & & & \multicolumn{5}{c}{OC20 S2EF Validation (\textbf{out-of-distribution (both)})} 
\\\cmidrule(r){4-7}
\multicolumn{1}{c}{}   &   &                   &      Energy MAE  &    Force MAE  & Force cos & EFwT \\
&   \multicolumn{1}{l}{Model}     &        &      \si{\milli\electronvolt} $\downarrow$   &    \si{\milli\electronvolt\per\angstrom} $\downarrow$     & $\uparrow$  & \% $\uparrow$ \\\cmidrule(r){1-7}
% & \multicolumn{1}{l}{Model} & &   &   \multicolumn{2}{c}{OC20-2M} & & \multicolumn{2}{c}{Other dataset} &  \\\cmidrule(r){5-6}\cmidrule(r){8-9}
% \multicolumn{9}{c}{} \\
\parbox[t]{2mm}{\multirow{6}{*}{\rotatebox[origin=c]{90}{\emph{same}}}} & Student (PaiNN-small)               &   &              $610$      &     $56.8$     &  $0.346$ &  $0.02$  \\
& Teacher (PaiNN-big)                  &   &              $554$        &   $59.2$     &   $0.379$      &   $0.03$    \\
& Vanilla KD (1)         &   &            $648 (\text{-}67.6\%)$          &          $61.1 (\text{-}179.2\%)$     &     $0.056 (\text{-}900.3\%)$        &  $0.02 (0\%)$   \\
& Vanilla KD (2)         &   &                $592 (32.2\%)$      &          $60.6 (\text{-}158.3\%)$     &     $0.310 (\text{-}112.4\%)$        &   $0.02 (0\%)$ \\
& n2n              &   &             $557 (94.2\%)$         &          $56.0 (\textbf{33}\%)$     &     $0.351 (14.8\%)$        & $0.018 (\text{-}15.8\%)$    \\
& v2v                &   &          $547 (\textbf{112.6\%})$            &          $56.5 (12.3\%)$     &     $0.352 (\textbf{17.3\%})$        & $0.025 (\textbf{43.3})$    \\
\multicolumn{7}{c}{} \\
\parbox[t]{2mm}{\multirow{5}{*}{\rotatebox[origin=c]{90}{\emph{similar}}}} & Student (SchNet)                  &   &              $1450$     &   $78.4$    &  $0.202$  &  $0$   \\
& Teacher (PaiNN-big)                  &   &              $554$        &   $59.2$     &   $0.379$      &   $0.03$    \\
& Vanilla KD (1)             &   &            $1350 (\textbf{11.2\%})$          &          $75.9 (13\%)$     &     $0.227 (\textbf{14.3\%})$        &  $0.0025 (\textbf{1.8\%})$ \\
& Vanilla KD (2)               &   &            $1358 (10.2\%)$          &          $75.7 (\textbf{14.1\%})$     &     $0.226 (13.8\%)$        &  $0 (0\%)$ &  \\
& n2n              &   &            $1396 (6.1\%)$          &          $76.2 (11.5\%)$     &     $0.220 (10.4\%)$        &  $0 (0\%)$  &  \\
\multicolumn{7}{c}{} \\
\parbox[t]{2mm}{\multirow{7}{*}{\rotatebox[origin=c]{90}{\emph{different}}}} & Student (PaiNN-big)                  &   &              $554$        &   $59.2$     &   $0.379$      &   $0.03$    \\
& Teacher (GemNet-OC)                      &   &              $370$        &   $31$     &   $0.606$      &   $0.23$     \\
& Vanilla KD (1)           &   &            $560 (\text{-}3.5\%)$          &          $53.7 (19.6\%)$     &     $0.373 (\text{-}2.3\%)$        &  $0.038 (3.1\%)$    \\
& Vanilla KD (2)             &   &            $511 (23.3\%)$          &          $323.1 (\text{-}935.7\%)$     &     $0.326 (\text{-}23.2\%)$        &  $0.027 (\text{-}2.6\%)$   \\
& n2n              &   &            $448 (\textbf{57.4\%})$          &          $52.8 (22.9\%)$     &     $0.384 (2.4\%)$        &  $0.058 (\textbf{13.1\%})$   \\
& e2n               &   &            $565 (\text{-}6.2\%)$          &          $50.9 (29.6\%)$     &     $0.394 (2.4\%)$        &  $0.046 (7.1\%)$     \\
& v2v              &   &            $537 (9.3\%)$          &          $50.8 (\textbf{29.7\%})$     &     $0.399 (\textbf{9\%})$        &  $0.036 (2.1\%)$        
% & e2e              &   &            $X (x\%)$          &          $X (x\%)$     &     $Y (y\%)$        &  $X (x\%)$  &          &       
\end{tabular}
\label{tab:ood_both_results}
\end{table}

\section{COLL experiments}
In Table~\ref{tab:coll_results}, we present the results on the COLL dataset \cite{gasteiger_dimenetpp_2020}, using different student-teacher configurations and KD strategies.
\label{ap:coll}
\begin{table}[H]
\centering
\caption{Evaluation results on the COLL test set.}
\begin{tabular}{llccccc}
\multicolumn{1}{c}{\multirow{3}{*}{}} & & \multicolumn{5}{c}{COLL test set} 
\\\cmidrule(r){4-7}
\multicolumn{1}{c}{}   &   &                   &      Energy MAE  &    Force MAE  & Force cos & EFwT \\
&   \multicolumn{1}{l}{Model}     &        &      \si{\milli\electronvolt} $\downarrow$   &    \si{\milli\electronvolt\per\angstrom} $\downarrow$     & $\uparrow$  & \% $\uparrow$ \\\cmidrule(r){1-7}
% & \multicolumn{1}{l}{Model} & &   &   \multicolumn{2}{c}{OC20-2M} & & \multicolumn{2}{c}{Other dataset} &  \\\cmidrule(r){5-6}\cmidrule(r){8-9}
% \multicolumn{9}{c}{} \\
\parbox[t]{2mm}{\multirow{6}{*}{\rotatebox[origin=c]{90}{\emph{same}}}} & Student (PaiNN-small)               &   &              $104.0$      &     $80.9$     &  $0.984$ &  $5.4$  \\
& Teacher (PaiNN-big)                  &   &              $85.8$        &   $64.1$     &   $0.988$      &   $10.1$    \\
& Vanilla KD (1)              &   &            $106.1 (\text{-}11.5\%)$          &          $82.0 (\text{-}6.5\%)$     &     $0.984 (2.3\%)$        &  $4.46 (\text{-}20.2\%)$    \\
& Vanilla KD (2)             &   &            $\textbf{86.4 (96.7\%)}$          &          $80.9 (0\%)$     &     $0.983 (\text{-}2.3\%)$        &  $4.3 (\text{-}23.7\%)$    \\
& n2n                &   &            $92.5 (63.2\%)$          &          $77.8 (18.5\%)$     &     $0.984 (18.2\%)$        &  $\textbf{6.63 (26.5\%)}$    \\
& v2v                       &   &            $90.4 (74.7\%)$   &          $\textbf{70.4 (62.5\%)}$     &     $\textbf{0.986 (45.5\%)}$        &  $5.8 (8.4\%)$    \\
\multicolumn{7}{c}{} \\
\parbox[t]{2mm}{\multirow{5}{*}{\rotatebox[origin=c]{90}{\emph{similar}}}} & Student (SchNet)                 &   &              $146.5$        &   $121.2$     &   $0.970$      &   $2.75$    \\
& Teacher (PaiNN-big)       &   &            $85.8$        &   $64.1$     &   $0.988$      &   $10.1$    \\
& Vanilla KD (1)            &   &            $146 (0.7\%)$          &         $121 (0.7\%)$     &     $0.970 (1.1\%)$        &  $2.54 (\text{-}2.9\%)$    \\
& Vanilla KD (2)            &   &            $\textbf{104 (69.9\%)}$          &          $121 (0.5\%)$     &     $0.970 (1.1\%)$        &  $\textbf{6.45 (50.7\%)}$    \\
& n2n                       &   &            $142 (8.1\%)$          &          $\textbf{117 (7.0\%)}$     &     $\textbf{0.971 (5.4\%)}$        &  $2.63 (\text{-}1.6\%)$    \\
\multicolumn{7}{c}{} \\
\parbox[t]{2mm}{\multirow{7}{*}{\rotatebox[origin=c]{90}{\emph{different}}}} & Student (PaiNN-big)                     &   &              $85.8$        &   $64.1$     &   $0.988$      &   $10.1$    \\
& Teacher (GemNet-OC)                  &   &              $44.8$        &   $38.2$     &   $0.994$      &   $20.2$    \\
& Vanilla KD (1)          &   &            $86.2 (\text{-}0.5\%)$          &          $63.9 (0.2\%)$     &     $0.988 (1.5\%)$        &  $10.1 (0.1\%)$    \\
& Vanilla KD (2)         &   &            $61.4 (28.5\%)$          &          $62.9 (1.9\%)$     &     $0.988 (5.2\%)$        &  $13 (29.7\%)$    \\
& n2n             &   &            $\textbf{60.4 (29.6\%)}$          &          $\textbf{61.2 (4.6\%)}$     &     $\textbf{0.989 (14.9\%)}$        &  $\textbf{13.6 (35.1\%)}$    \\
& e2n                       &   &            $77.3 (9.9\%)$          &          $63.3 (1.2\%)$     &     $0.988 (7.9\%)$        &  $11 (9.4\%)$    \\
& v2v                  &   &            $81.2 (5.4\%)$          &          $63.3 (1.2\%)$     &     $0.988 (3.4\%)$        &  $10.5 (4.7\%)$    \\   

\end{tabular}
\label{tab:coll_results}
\end{table}

\section{Data augmentation}
\label{ap:data_aug}

\subsection{Data jittering} 
\label{ap:jittering}
To create additional data, one approach we tried was to add noise to the atomic positions of the training samples and then use the teacher to label the samples. Apart from adding random noise, we also tried to optimize the positions using gradient ascent such that the difference between the predictions of the student and teacher were maximized. Denoting the noise as $\delta$, we obtain the noise atomic positions as $\boldsymbol{X}_\delta = \boldsymbol{X} + \delta$. Let the student and teacher models be denoted as $f_\text{s}$ and $f_\text{t}$ respectively, we then obtained the noise $\delta'$ by initializing this as $\boldsymbol{0}$ and
\begin{align*}
            \loss_{\text{KD}} &= \loss_0(f_s(\boldsymbol{X}_{\delta}, \boldsymbol{z}), f_t(\boldsymbol{X}_{\delta}, \boldsymbol{z})) \\
            \delta' &= \delta + \alpha \nabla_{\delta} \mathcal{L}_{\text{KD}}.
\end{align*}

This becomes computationally expensive, though, as it requires additional gradients for $\nabla_{\delta} \mathcal{L}_{\text{KD}}$. Also, when comparing this with using random noise, we did not see any improvements.

\subsection{Synthetic data} 
\label{ap:synthetic}

\paragraph{Combined dataset 2M+d1M.}
We generated 1M synthetic samples by first drawing 100k random adsorbate and catalyst combinations (systems) and then running relaxations with a pre-trained GemNet-OC model. Out of these relaxations with 100 steps on average (200 max), we randomly draw approx. 10\% to obtain 1M samples.

In the next step, we combine the 1M samples with the 2M OCP dataset, which is based on DFT relaxations. Directly working with this combined dataset means iterating over a 1-to-2 ratio of samples from each subset in an epoch of 3M samples. To control this ratio, we define the target ratio of samples from the synthetic dataset during training $\alpha_{target}\in[0,1]$. Setting  $\alpha_{target}=0.5$ means that per epoch we iterate over the 1M dataset 1.5 times and over the 2M dataset 0.75 times.

\paragraph{Different weighting in loss depending on origin.}

Next to specifying a sampling ratio of samples from the synthetic dataset, we can also specify how to weight the contribution of samples to the loss based on their origin (DFT or synthetic). To achieve this, we specify the weighting ratio of synthetic to DFT samples $r_{\text{s/dft}} = w_{\text{s}}/w_{\text{dft}} \in \mathbb{R}^+$. In each batch, we compute a weighting factor in front of the synthetic $w_\text{s}$ and DFT $w_{\text{dft}}$ samples satisfying the conditions 
\begin{equation}
    w_\text{s} \cdot  \alpha_{\text{batch}} + w_{\text{dft}} \cdot (1-\alpha_{\text{batch}}) = 1,
\end{equation}
where $\alpha_{\text{batch}}$ is the ratio of synthetic to DFT samples in a batch. Hence, when we combine DFT and synthetic data - $\alpha_{\text{batch}} \in(0,1)$, we derive the following weights:
\begin{align}
  w_{\text{dft}}&=(1-\alpha_{\text{batch}}+\alpha_{\text{batch}} \cdot r_{\text{s/dft}})^{-1},\\
  w_{\text{s}}&=r_{\text{s/dft}} \cdot w_{\text{dft}}.
\end{align}
Likewise, when $\alpha_{\text{batch}}=\{0, 1\}$ - i.e., we either train on DFT or synthetic data exclusively, the corresponding weighting coefficient ($w_\text{s}$ or $w_{\text{dft}}$) is naturally equal to $1$.

% Samples in the OC20 datasets are generated by first running DFT-based relaxations on randomly sampled systems, i.e. combination of adsorbate and catalyst, and subsequently randomly sampling individual frames along these relaxation trajectories. By sampling along the same trajectory, many samples are correlated, which limits the diversity of the dataset. To tackle this issue and explore the capabilities of knowledge distillation, we generated our own distilled 1M (d1M) dataset by sampling novel systems (generated with the OC Datasets codebase \footnote{\url{https://github.com/Open-Catalyst-Project/Open-Catalyst-Dataset}}) and running relaxations with a pre-trained GemNet-OC 2M model \footnote{\url{https://github.com/Open-Catalyst-Project/ocp/blob/main/MODELS.md}}.

% We tried different protocols for incorporating this new d1M dataset, all of which are based on jointly training on the 2M as well as the d1M datasets. To explore different dataset combinations, we introduce two new hyperparameters: the ratio of samples drawn from the d1M dataset $\alpha_{d1M}$, and the loss weighting ratio between DFT- and GemNet-OC-based samples $r_{d1M/DFT}$. However, after extensively exploring these two hyperparameters for training our PaiNN-big model on the 2M+d1M data, we could not see improvement over the baseline performance. This result is surprising as similar approaches have been used for knowledge distillation in other domains \cite{xu2022contrastive}, and we hypothesize that by the out-of-distribution nature of the systems in the d1M dataset, some of the labels from the pre-trained GemNet-OC model might lead the student's training in the wrong direction. For completeness, we also explored knowledge distillation on a combination of both datasets on the n2n as well as v2v distillation from GemNet-OC to PaiNN-big, but didn't observe significant improvement either.



\section{Hyperparameter studies}
\label{ap:hyperparams}
We additionally investigated different aspects of our KD protocols, which we present below. The results in this section are on a smaller validation set, containing 30 thousand randomly sampled systems from the in-distribution validation set.

\subsection{Effect of transformations} 
\label{ap:transformations}
We first explore the effect of the transformations $\mathcal{M}_\text{s}, \mathcal{M}_\text{t}$ on the performance of KD. To reduce the search space, we fix $\mathcal{M}_\text{t}$ as the identity transformation and only explore the choice of $\mathcal{M}_\text{s}$. We conduct our experiments in the setting of distilling knowledge from GemNet-OC to PaiNN with our \emph{n2n} and \emph{e2n} KD strategies, and compare the following functions: the identity transformation; a linear transformation; and a more powerful two-layer MLP. Note that we cannot include the identity transformation for \emph{n2n} as the node features in the respective models are of different dimensions. For \emph{n2n}, we also include a comparison with the Global Structure Preserving (GSP) loss \cite{joshi2022}. As this does not rely on the features having the same dimensions, we try using both an identity transformation and a learned linear layer of the feature before computing the GSP loss. 

From the results, presented in Table~\ref{tab:projhead}, we conclude that a simple linear layer achieves the best KD performance and is enough to transform the features of the student model.

\begin{table}[h]
\centering
\caption{\label{tab:projhead} Investigating the effect on performance of different functions used as $\mathcal{M}_\text{s}$ transformations of student features. GPS is Global Structure Preserving \cite{joshi2022}. }
\begin{tabular}{cccc}
                         &                 & \multicolumn{2}{c}{Relative Improvement (\%)} \\ \cline{3-4} 
Loss                     & Transformation  & Energy MAE            & Force MAE            \\ \hline
\multicolumn{1}{c|}{}    & Identity        & -                     & -                    \\
\multicolumn{1}{c|}{}    & Linear          & 0,0                   & 0,0                  \\
\multicolumn{1}{c|}{n2n} & Projection head & -10,0                 & -6,3                 \\
\multicolumn{1}{c|}{}    & Identity + GSP  & -36,4                 & -7,2                 \\
\multicolumn{1}{c|}{}    & Linear + GSP    & -33,6                 & -10,6                \\ \hline
\multicolumn{1}{c|}{}    & Identity        & 0,0                   & 0,0                  \\
\multicolumn{1}{c|}{e2n} & Linear          & -0,4                  & 0,2                  \\
\multicolumn{1}{c|}{}    & Projection head & -9,1                  & -7,0                
\end{tabular}
\end{table}


\subsection{Effect of losses}
\label{ap:losses}
We then explored the effect of the $\loss_\text{feat}$ loss. We conducted our experiments in the settings of doing \emph{n2n} KD on PaiNN with GemNet-OC as a teacher. We investigated the following losses $\loss_\text{feat}$: MSE; MAE; and directly optimize layer similarity using CKA as a distillation loss. The results are in Table~\ref{tab:mae}. We do not present any results for CKA, as it showed poor results from the beginning, with forces worse than training without any KD loss at all, and was, therefore, not further trained.
%included for further hyperparameter tuning (i.e. optimization of the $\lambda$ weight factor of the KD loss$\loss_\text{KD}$). 
The results (\cref{tab:mae}) demonstrated that the difference between MSE and MAE is marginal.

\begin{table}[h]
\centering
\caption{\label{tab:mae} Investigating the effect on performance of different loss functions used as a $\loss_\text{feat}$ loss. We also explored CKA similarity as loss, but this showed poor results and was not included for further hyperparameter tuning. }
\begin{tabular}{ccc}
               & \multicolumn{2}{c}{Relative Improvement (\%)} \\ \cline{2-3} 
Loss           & Energy MAE            & Force MAE            \\ \hline
MSE (baseline) & 0.0                   & 0.0                  \\
MAE            & -1.7                  & 1.2                  \\
CKA            & -                     & -                   
\end{tabular}
\end{table}

\subsection{Effect of feature selection}
\label{ap:features}
% \textit{We complemented our experiments by performing CKA-based similarity analyses that allowed us to gain insights into how these protocols change the architecture of the student.}
Inspired by general KD setups, our initial experiments with GemNet-OC and PaiNN were based on the distillation of features across the final layers of the two models. However, we decided to utilize a separate approach, where we distil information from all layers of the teacher instead of the final one only. To achieve that, we exploited the fact that, before making predictions, the GemNet-OC model concatenates the output features across all layers, passes them through an MLP-based fusion function and derives final aggregated node- and edge-level output features. Hence, we decided to use these final aggregated features as a means of distilling information from all layers of the teacher. The results in the settings of \emph{n2n} (\cref{tab:features}) showed a drastic improvement in performance compared to using the features from the usual final layer of the teacher. Hence, we used that approach throughout all experiments with GemNet-OC and PaiNN as teacher and student respectively. We also explored the utility of combining more features from the student (\cref{tab:features}), but this did not seem to be of substantial benefit.

\begin{table}[h]
\centering
\caption{\label{tab:features} Investigating the effect on performance of feature selection in feature-based KD. We explore different choices of node features to distil when doing \emph{n2n} KD from GemNet-OC into PaiNN. CONCAT+MLP is the resulting aggregated features after concatenating the output blocks and applying an MLP, i.e., the final features used for making predictions in GemNet-OC. H4 is the feature after the final interaction block in GemNet-OC, before the last output block. U6 and M6 are the node features after the sixth and final message and update blocks in PaiNN, respectively.}
\begin{tabular}{cccc}
                &                              & \multicolumn{2}{c}{Relative Improvement (\%)} \\ \cline{3-4} 
GemNet-OC layer & PaiNN layer                  & Energy MAE             & Force MAE            \\ \hline
CONCAT+MLP      & \multicolumn{1}{c|}{U6}      & 0.0                    & 0.0                  \\
H4              & \multicolumn{1}{c|}{U6}      & -31.4                  & -4.4                 \\
CONCAT+MLP      & \multicolumn{1}{c|}{M6 + U6} & 0.2                    & 0.5                    
\end{tabular}
\end{table}

\paragraph{Additional ablation studies} 
To make our insights into the importance of feature selection, we performed a series of ablation studies, where we varied the feature layer we utilized for KD in both models, as well as the number of layers we use from the student model. The results of these analyses are presented in
\cref{tab:ablation1} - \cref{tab:ablation3}.

It is noteworthy to indicate that due to the number of experiments, if not mentioned otherwise, the models presented in these 3 tables have not been subject to full hyperparameter tuning to find the best performing $\lambda$ weighting KD factor. To reduce the time burden, we have instead conducted a faster optimization, where $\lambda$ was selected based on the consistency of the resulting total loss (in the initial stages of the training run) with that in previously tuned experiments. More specifically, for each setting, we selected the $\lambda$ value that results in a distillation loss $\lambda \loss_{\text{KD}}$ that, at the beginning of the training process, was the most similar in magnitude to the one achieved with our default KD setup (after tuning) - i.e., when using the final feature in PaiNN and the concatenated output block features in GemNet-OC.
 
\begin{table}[h]
\caption{Ablation study across the student architecture. The table summarizes the performance of the student model after distilling knowledge from the final layer/block of the teacher (GemNet-OC with its U6 feature) into different blocks/layers of the student (PaiNN). M and U indicate whether the feature taken is after a message or an update layer, whereas the number indicates which block, in particular, was used - i.e., U2 refers to the output of the second update block). We observe a distinct trend that performance is hurt as we go deeper into the student architecture.}
\label{tab:ablation1}
\centering
\begin{tabular}{ccc}
                                   & \multicolumn{2}{c}{Relative Difference (\%)} \\ \cline{2-3} 
Student Layer                              & Energy MAE            & Force MAE            \\ \hline
\multicolumn{1}{c|}{M1}            & -38,6                 & -4,9                 \\
\multicolumn{1}{c|}{U1}            & -37,2                 & -4,1                 \\
\multicolumn{1}{c|}{M2}            & -37,1                 & -4,1                 \\
\multicolumn{1}{c|}{U2}            & -34,1                 & -4,5                 \\
\multicolumn{1}{c|}{M3}            & -32,7                 & -4,3                 \\
\multicolumn{1}{c|}{U3}            & -28,5                 & -5,5                 \\
\multicolumn{1}{c|}{M4}            & -14,7                 & 2,6                  \\
\multicolumn{1}{c|}{U4}            & -11,2                 & 2,5                  \\
\multicolumn{1}{c|}{M5}            & -9,6                  & 3,6                  \\
\multicolumn{1}{c|}{U5}            & -7,9                  & 3,3                  \\
\multicolumn{1}{c|}{M6}            & -7,2                  & 3,0                  \\
\multicolumn{1}{c|}{U6 (baseline)} & 0,0                   & 0,0                 
\end{tabular}
\end{table}


\begin{table}[h]
\centering
\caption{Ablation study across the teacher architecture. The table summarizes the performance of the student model after distilling knowledge from different layers/blocks of the teacher (GemNet-OC) into the final layer/block of the student (PaiNN with its CONCAT + MLP feature). H refers to the standard node features in GemNet-OC, whereas X indicates that it is the feature from an output block. The number indicates after which block in the architecture was the feature extracted. CONCAT+MLP is the feature after concatenating the output block features and processing this with an MLP. Here, we observe a similar pattern when considering energy predictions. For force predictions, however, we notice that distilling information from the output blocks can be beneficial, whereas all node features result in a drop in performance with respect to forces, too.}
\label{tab:ablation2}
\begin{tabular}{ccc}
\multicolumn{1}{l}{}                         & \multicolumn{2}{c}{Relative Difference (\%)} \\ \cline{2-3} 
Teacher Layer                                        & Energy MAE            & Force MAE            \\ \hline
\multicolumn{1}{c|}{H0}                      & -46.8                 & -10.6                \\
\multicolumn{1}{c|}{X0}                      & -23.4                 & 2.9                  \\
\multicolumn{1}{c|}{H1}                      & -55.2                 & -11.6                \\
\multicolumn{1}{c|}{X1}                      & -37.9                 & -3.8                 \\
\multicolumn{1}{c|}{H2}                      & -37.5                 & -9.7                 \\
\multicolumn{1}{c|}{X2}                      & -15.3                 & 3.0                  \\
\multicolumn{1}{c|}{H3}                      & -25.7                 & -8.3                 \\
\multicolumn{1}{c|}{X3}                      & -16.3                 & 3.1                  \\
\multicolumn{1}{c|}{H4}                      & -31.4                  & -4.4                \\
\multicolumn{1}{c|}{X4}                      & -16.0                 & 2.9                  \\
\multicolumn{1}{c|}{CONCAT + MLP (baseline)} & 0.0                   & 0.0                 
\end{tabular}
\end{table}


\begin{table}[h]
\centering
\caption{Ablation study across combinations of layers across the student architecture. The table summarizes the performance of the student model after distilling knowledge from the final layer/block of the teacher (GemNet-OC with its CONCAT + MLP feature) into different combinations of layers/blocks within the student (PaiNN). A given combination is used by concatenating the respective features and passing them through a linear layer that is learnt during training. Features denoted as in \cref{tab:ablation1}. We did not observe any distinctive trends or performance gains, and hence decided to use the final U6 in PaiNN only in further experiments.}
\label{tab:ablation3}
\begin{tabular}{ccc}
\multicolumn{1}{l}{}               & \multicolumn{2}{c}{Relative Difference (\%)} \\ \cline{2-3} 
Student Layer(s)                           & Energy MAE            & Force MAE            \\ \hline
\multicolumn{1}{c|}{U6 (baseline)} & 0.0                   & 0.0                  \\
\multicolumn{1}{c|}{U5 + U6}       & -2.6                  & 2.0                  \\
\multicolumn{1}{c|}{U4 +U5 + U6}   & -10.9                 & 4.3                  \\
\multicolumn{1}{c|}{M6 + U6}       & 0.5                   & 0.5                  \\
\multicolumn{1}{c|}{M5 + U5 + U6}  & -2.4                  & 2.0                 
\end{tabular}
\end{table}


% \subsection{Effect of combining strategies}
% \label{ap:combinations}
% ...

\section{Explainability}
\label{ap:cka}

% \begin{figure}
%     \centering
%     \includegraphics[width=0.9\textwidth]{cka/after_n2n.png}
%     \includegraphics[width=0.9\textwidth]{cka/after_n2n.png}
%     \includegraphics[width=0.9\textwidth]{cka/after_n2n.png}
%     \caption{Other features - GemNet-OC->PaiNN: no KD (n2n, e2n, v2v features).} 
%     \label{fig:cka_other_features}
% \end{figure}

We utilize CKA similarity scores to monitor the effect of KD throughout our studies. Here, we present a selection of the analyses we have performed, summarizing how KD influences the similarity between teacher and student models across teacher-student configurations (\cref{fig:cka_after_kd_all_models}); across KD protocols (\cref{fig:cka_after_different_kd}); and feature selections  (\cref{fig:cka_after_n2n_different_features}). We found out that such similarity metrics can be effectively used to examine and profile different KD approaches, as well as as a potential debugging tool. We also explored the utility of CKA (in conjunction with measures of the predictive ability of individual features) to inform the design of (optimal) KD strategies and feature selection protocols \textit{a priori}, but the results were not conclusive to include here.


\begin{figure}[h]
    \centering
    \includegraphics[width=\textwidth]{cka/different_models.png}
    \caption{We explore the effect of \emph{n2n} KD on the feature similarity between different student-teacher configurations: (a) GemNet-OC -> PaiNN; (b) PaiNN-big -> PaiNN-small; (c) PaiNN-big -> SchNet. The layer pair that was used in each experiment is indicated with a  $\textcolor{cyan_our}{\swarrow}$. Note the scale difference in the \textit{Similarity gain} plots.} 
    \label{fig:cka_after_kd_all_models}
\end{figure}

\begin{figure}[h]
    \centering
    \includegraphics[width=\textwidth]{cka/different_kd.png}
    \caption{We explore the effect of different KD strategies - \emph{n2n} and \emph{e2n} KD on the feature similarity between the student and the teacher models. This is computed for GemNet-OC -> PaiNN: (a) \emph{n2n}; (b) \emph{e2n}. The layer pair that was used in each experiment is indicated with a $\textcolor{cyan_our}{\swarrow}$. Note the scale difference in the \textit{Similarity gain} plots.} 
    \label{fig:cka_after_different_kd}
\end{figure}

\begin{figure}[h]
    \centering
    \includegraphics[width=\textwidth]{cka/feature_selection.png}
    \caption{We explore the effect of feature selection in KD on feature similarity between the student and the teacher models: (a) H4->U6; (b) CONCAT+MLP->U6; (c) CONCAT+MLP->M6+U6; (d) H4->U5; (e) X2->M2. The layer pair that was used in each experiment is indicated with a $\textcolor{cyan_our}{\swarrow}$. Note the scale difference in the \textit{Similarity gain} plots.} 
    \label{fig:cka_after_n2n_different_features}
\end{figure}

\end{document}